\documentclass[11pt]{article}
\usepackage{fontspec}
\setmainfont{Times New Roman}
\setsansfont{Arial}
\setmonofont{Consolas}
\usepackage{amsmath,amssymb,mathtools,array,booktabs}
\usepackage{geometry}
\usepackage{microtype}
\geometry{a4paper,margin=2.35cm}
\setlength{\parindent}{1.25em}
\setlength{\parskip}{0pt}
\makeatletter
\renewcommand\footnotesize{\@setfontsize\footnotesize{7.7}{8.7}\selectfont}
\makeatother
\emergencystretch=35em
\allowdisplaybreaks
\sloppy
\hbadness=10000
\vbadness=10000
\hfuzz=2pt
\vfuzz=10pt
\raggedbottom
\providecommand{\Amod}{\mathfrak A}
\providecommand{\Bmod}{\mathfrak B}
\providecommand{\Gmod}{\mathfrak G}
\providecommand{\Mmod}{\mathfrak M}
\providecommand{\Nmod}{\mathfrak N}
\providecommand{\Rmod}{\mathfrak R}
\providecommand{\Smod}{\mathfrak S}
\providecommand{\mideal}{\mathfrak m}
\providecommand{\nideal}{\mathfrak n}
\providecommand{\gideal}{\mathfrak g}
\providecommand{\hideal}{\mathfrak h}
\providecommand{\jideal}{\mathfrak j}
\providecommand{\tideal}{\mathfrak t}
\providecommand{\aideal}{\mathfrak a}
\providecommand{\bideal}{\mathfrak b}
\providecommand{\bb}{\mathfrak b}
\providecommand{\cc}{\mathfrak c}
\providecommand{\zero}[1]{\equiv 0\pmod{#1}}
\providecommand{\Norm}{N}

\begin{document}

% Унит: Папер22 интродуцтион воркинг Интерславиц Латин пацкет. Дирецт Интерславиц воркинг транслатион фром те Герман финал-аудитед соурце слице; Латин ис те ауториты лане; нот а финал цхецкпоинт.
\setcounter{footnote}{0}
\section*{22. Обработка К. Хентзелт:\\ к теорији полиномных идеалов и резултантов}
\begin{center}
Math. Анн. 88 (1923), стр. 53--79
\end{center}

В следујучем иде о \emph{обчу задачу теорије елиминације: определити обче нулы всих полиномов идеала из полиномов}; або, што јест с тым идентично, нулы либовольного конечного числа полиномов, кторе творјут базис идеала; при том једночасно појави се поньатијно толкованье наставајучих многократностиј. Коефициенты полиномов можно прирадити либовольному польу; нулы тогда належат алгебраично замкньеному польу, изведеному из ньега. Кроме того идеал можно без ограниченьа обчности сматрати трансформованым (§~3). Тогда сушчствены резултат јест следујучи:

\emph{Каждому идеалу $\mideal$ из полиномов можно прирадити резултантну форму:}
\[
 R_{\mideal}=R^{(1)}(x_1,\ldots,x_n)\cdots
 R^{(i)}(x_i,\ldots,x_n)\cdots R^{(n)}(x_n)\zero{\mideal},
\]
\emph{тако же $R_{\mideal}$ изчезаје за все нулы $\mideal$ и само за ньих. Ако $\nideal$ јест делитель од $\mideal$, и резултантне формы $R_{\nideal}$ и $R_{\mideal}$ согласујут се, тогда согласујут се такого идеалы $\nideal$ и $\mideal$.}\footnote{Курт Хентзелт од октобра 1914 јест безвестно изчезлы при Диксмуиден и мора быти числьены меджу умрлими. Тут дана разправа јест полно свободна обработка најсушчственеј чести его дисертације ``Zur Theorie der Форменмодулн und Ресултантен'', с ктороју он в лету 1914 докторовал при Е. Фисцхер в Ерланген. Та дисертација, написана целком на основе его властных идеј, јест без пробојев изграджена; але лема следује лему, все поньатја сут описане формулами с четырми и петими индексами, текст скоро полно отсутствује, тако же разуменьу ставјут се највечше трудности. До плановане преработкы он сам уже не дошел. Ја подају работу в чисто поньатијној форме, чим достигаје се велико упрошченье доказов, кторы в основных мысльах повсуду идут назад к Хентзелт, и, како надејам се, очивидна става красота работы. -- Чести дисертације, кторе -- при даном базису -- односет се к питаньу творјеньа наставајучих функциј конечным числом кроков, имајут остати резервоване за одделно опубликованье. (Е. Н.)}

Друга чест главного теорема показује, же резултантна форма дава нулы в характерној многократности. Та друга чест јест аналогон тому, же два идеала алгебраичного числового польа, од кторых једен јест делитель другого, согласујут се тогда и само тогда, когда ньих \emph{нормы} согласујут се; внутрны основ обојих теоремов јест такого тен самы.

Именно, $\mideal$ можно поймати како модул линеарных форм $\mathfrak M_{i-1}$, ако кажды полином поймати како линеарну форму в степенных продуктах од $x_1,\ldots,x_{i-1}$ и присојединити $x_{i+1},\ldots,x_n$ к области коефициентов. Резултантны фактор $R^{(i)}$ тогда става равны \emph{норме} одповедајучего основного модула (§~1) $\Gmod_{i-1}$ по $\mathfrak M_{i-1}$; из того такој даје се \emph{толкованье многократности -- продукт степена и експонента фактора $R^{(i)}$ -- чрез число линеарно независных класов остатков}\footnote{Ти последни резултаты показывајут своје право значенье, когда привлече се разложенье идеалов на примарне; многократност, одповедајуча примарному идеалу, тогда определьа се како число линеарно независных класов остатков комплемента по том идеалу; о том буде говорјено на другом месте. (Е. Н.)}, факт, кторы доселе был знаны само в специалном случају, кде резултант можно определити за неопредельене коефициенты\footnote{Ср. Мацаулаы, \emph{The algebraic theory of modular systems}, Цамбридге Трацтс 19 (Цамбридге Университы Пресс, 1916), Но. 67; такого Ласкер, Zur Theorie der Модулн und Ideale, Math. Анн. 60 (1905), стр. 20--116; стр. 98. Же в том случају резултант и резултантна форма согласујут се, показано јест в еще не опубликованых теоремах Хентзелт, споменутых под 1). (Е. Н.)}.

За выводженье тих резултатов в §~1 и §~2 дајут се теоремы о модулах из линеарных форм, кторых коефициенты сут целе рационалне числа або полиномы једнеј променној. Связ с идеалами из полиномов, уведеными в §~3, посредкује теорем изоморфизма §~4. §~5 дава выше споменуту другу чест главного теорема, §~7 теорем о нулах, кторему в §~6 претоди обча теория елиминације. Тамо се једночасно, при уведеньу новых променных, доказује разложенье резолвенты на линеарне формы тих променных; и тым при тут даној елиминацији дава се безпорочны доказ за једно тврдјенье Кронецкер\footnote{Покус доказа у König, \emph{Algebraische Größen} (Леипзиг 1903, Teubner), V, §~4 -- за Кронецкерову теорију елиминације -- како знано, не успел. Же и при многократностјах, кторе König уведл в Кронецкерову теорију елиминације, друга чест Хентзелтового главного теорема не важи, показује Мацаулаы, лоц. цит., стр. 23; там дале показано јест, же Кронецкерова теория елиминације не одповедаје разложеньу на примарне идеалы. (Е. Н.)}.

\section*{§ 1. Модулы из линеарных форм.}

Под целыми величинами $a,b,c,\ldots$ в првих двох параграфах разумејут се целе рационалне числа или полиномы једнеј неопредельенеј с коефициентами из либовольного, абстрактно определьеного польа $P$; под линеарными формами $a(\xi),b(\xi),\ldots$ разумејут се таке формы в конечном числу неопредельеных $\xi_1,\ldots,\xi_k$ с целыми величинами како коефициентами.

\emph{Модул $\Amod$ из линеарных форм} определьа се условијами: заедно с $a(\xi)$ и $b(\xi)$ модул $\Amod$ содержи такоже разлику $a(\xi)-b(\xi)$; заедно с $a(\xi)$ он содержи такоже $c\cdot a(\xi)$, кде под $c$ разумеје се либовольна цела величина.

Под \emph{рангом $\Amod$} разумејемо, како обычно, максимално число линеарно независных линеарных форм из $\Amod$; под $\lambda$-тым \emph{детерминантным делительем} $d_\lambda$ модула $\Amod$ -- највечши обчи делитель всех детерминантов реда $\lambda$ из $\Amod$ (то јест всех детерминантов реда $\lambda$, кторе можно утворити из матрице коефициентов либоволных $\lambda$ линеарных форм из $\Amod$). Ако $\rho$ јест ранг $\Amod$, тако же $d_1,\ldots,d_\rho$ сут ненулове детерминантне делительи, тогда \emph{елементарне делительи} модула $\Amod$ определьајут се формулами $e_1=d_1$, $e_2=d_2/d_1,\ldots,e_\rho=d_\rho/d_{\rho-1}$, $e_{\rho+i}=0$. Како знано, $\Amod$ јест конечны модул, кторы имаје модулны базис точно из $\rho$ линеарных форм $a_1(\xi),\ldots,a_\rho(\xi)$; чрез ньих всака линеарна форма из $\Amod$ може се представити линеарно, с целыми величинами како коефициентами:
$\Amod=(a_1(\xi),\ldots,a_\rho(\xi))$. Ако $A$ означује матрицу коефициентов тих линеарных форм, тогда детерминантне и елементарне делительи $\Amod$ согласујут се с одповедајучими делительами $A$; из того најпрво следује, же \emph{всакы элементарны делитель $e_i$ дели следујучи $e_{i+1}$}. Матрично равнанье, дано теоријеју елементарных делительев,
\[
 P A Q=\begin{pmatrix}
 e_1&&&0\\
 &e_2&&\\
 &&\ddots&\\
 0&&&e_\rho
 \end{pmatrix}
\]
дале показује -- ако чрез унимодуларну трансформацију $\xi=Q(\eta)$ введут се нове неопредельене $\eta_1,\ldots,\eta_k$ -- же обстаје модулно равнанье:
\begin{equation}
 \Amod=(e_1\eta_1,e_2\eta_2,\ldots,e_\rho\eta_\rho).
\tag{1}
\end{equation}

Најсушчественејши појем за то, что следује, јест појем основного модула\footnote{Ср. Steinitz, Рецхтецкиге Сыстеме und Модулн in algebraischen Zahlkörpern II. Math. Анн. 72 (1912), стр. 297--345, Но. 36. Хентзелт, како зда се, во всаком случају независно од Стеинитза -- с кторым и индје најдут се точкы дотыка -- дошел јест за свой простейши случай к тому појму в облику формулы (4). (Е. Н.)}, по следујучеј

\medskip
\noindent\textbf{Дефиниция I.} \emph{Модул $\Gmod$ называје се основным модулом, ако он не имаје властного делительа једнакого ранга}\footnote{Делитель разумены в модулном смыслу: $\Amod$ јест делимы чрез $\Bmod$, $\Amod\zero{\Bmod}$, ако всакы элемент из $\Amod$ содержи се в $\Bmod$; $\Bmod$ называје се властным делительем, ако он содержи елементы различне од елементов $\Amod$.}.

Тогда $\Gmod$ јест основны модул тогда и само тогда, когда из
\begin{equation}
 c\cdot g(\xi)\zero{\Gmod};\quad c\ne0 \quad\text{всагда следује:}\quad g(\xi)\zero{\Gmod}.
\tag{2}
\end{equation}

Связ меджу либовольным модулом и основным модулом даје следујучи

\medskip
\noindent\textbf{Теорем I.} \emph{Всакы модул $\Amod$ јест делимы чрез једин и само једин основны модул $\Gmod$ једнакого ранга; $\Gmod$ определьа се како найменши делитель $\Amod$ једнакого ранга и представјаје се чрез}
\begin{equation}
 \Gmod=(\eta_1,\ldots,\eta_\rho);
\tag{3}
\end{equation}
\emph{$\Gmod$ буде называны основным модулом модула $\Amod$.}

Услов’е, же $\Gmod$ ма быти найменши делитель једнакого ранга, выража се тако: $\Gmod$ содержи все и само те линеарне формы $g(\xi)$, кторе задовольајут релацији
\begin{equation}
 b\cdot g(\xi)\zero{\Amod};\quad b\ne0.
\tag{4}
\end{equation}
Из того најпрво следује, же $\Gmod$ јест модул, бо заедно с
\[
 b_1g_1(\xi)\zero{\Amod};\quad b_1\ne0;
 \qquad
 b_2g_2(\xi)\zero{\Amod};\quad b_2\ne0
\]
изполньено јест такоже
\[
 b_1b_2\bigl(g_1(\xi)-g_2(\xi)\bigr)\zero{\Amod};\quad b_1b_2\ne0,
\]
и, разуме се, такоже $b_1\cdot c g_1(\xi)\zero{\Amod}$. Але $\Gmod$ јест такоже основны модул; бо из
\[
 c\cdot g(\xi)\zero{\Gmod};\quad c\ne0
\]
по (4) следује
\[
 b c g(\xi)\zero{\Amod};\quad bc\ne0,
\]
а зато, знову по (4), такоже $g(\xi)\zero{\Gmod}$.

Дале не обстаје основны модул $\Gmod_1$, различны од најменшего делительа $\Gmod$ једнакого ранга, кторы бы задовольавал те условија. Бо $\Gmod_1$ мусил бы быти делимы чрез $\Gmod$, але како основны модул он не имаје властного делительа једнакого ранга, и зато јест идентичны с $\Gmod$. Наконец модул представјены правоју страноју (3) јест основны модул, бо всакы властны делитель мусил бы содержати дальньу неопредельену $\eta$, значи имал бы вышши ранг. Тымы (3) даје еднозначно определьены основны модул од $\Amod$, и \emph{Теорем I јест доказаны во всех честјах}.

Дальны појем, сушчествены за то, что следује, јест Дедекиндов појем квоциента двох модулов\footnote{При Дедекинду иде о числове модулы, за кторе определьено јест такоже множенье, тако же Дедекиндов квоциент не согласује се с квоциентом определьеным тут. (Е. Н.)}, именно:

\medskip
\noindent\textbf{Дефиниция II.} \emph{Квоциент $\cc=\Amod/\Bmod$ двох модулов из линеарных форм определьа се како целост целых величин $c$, кторе задовольајут условју}
\[
 c\cdot\Bmod\zero{\Amod}.
\]
\emph{$\cc$ представльа идеал из целых величин, зато главны идеал. Одповедајуче, квоциент $\mathfrak C=\Amod/\bb$ модула $\Amod$ из линеарных форм чрез идеал $\bb$ из целых величин определьа се како целост линеарных форм $c(\xi)$, кторе задовольајут условју}
\[
 c(\xi)\cdot\bb\zero{\Amod}.
\]
\emph{$\mathfrak C$ знову представльа модул из линеарных форм, именно, скоро $\bb$ јест различны од нулового идеала, модул једнакого ранга како $\Amod$.}

Тут треба само приметити, же из дефиниције ясно јест: квоциент всагда имаје модулно свойство; системы из целых величин с модулным својством сут пак идеалы.

Појем квоциента води к связи меджу основным модулом и највышшим елементарным делительем:

\medskip
\noindent\textbf{Теорем II.} \emph{Меджу основным модулом и највышшим елементарным делительем од $\Amod$ обстаје взаимна релација}
\begin{equation}
 (e_\rho)=\Amod/\Gmod;\qquad \Gmod=\Amod/(e_\rho),
\tag{5}
\end{equation}
\emph{кде $(e_\rho)$ значи главны идеал изведены из $e_\rho$.}

Бо понеже всакы элементарны делитель дели $e_\rho$, из (1) и (3) добивајемо
\begin{equation}
 e_\rho\Gmod\zero{\Amod}\quad \text{и следовательно}\quad (e_\rho)\cdot\Gmod\zero{\Amod};
\tag{6}
\end{equation}
зато $(e_\rho)$ јест делимы чрез квоциент $\Amod/\Gmod$. Але важи и обратно; бо из
\[
 b\Gmod\zero{\Amod}
 \quad\text{следује}\quad b\eta_\rho\zero{\Amod}
 \quad\text{и зато}\quad b\zero{(e_\rho)},
\]
чим доказана јест прва формула (5)\footnote{Ако положимо $\Amod_0=\Amod,\ldots,\Amod_\lambda=(\Amod,\eta_\rho,\ldots,\eta_{\rho-\lambda+1})$, кде всакы $\Amod_i$ имаје највышши элементарны делитель $e_{\rho-i}$, тогда по тыцх самых закльучаньах приходит:
\[
 (e_\rho)=\Amod_0/\Amod_1,\ldots,(e_{\rho-\lambda})=\Amod_\lambda/\Amod_{\lambda+1},\ldots,(e_1)=\Amod_{\rho-1}/\Amod_\rho.
\]
Обчејше, нехай положено буде $\zeta_\lambda=b_{\lambda1}\eta_1+\cdots+b_{\lambda\lambda}\eta_\lambda$, нехай $(b_{\lambda\lambda},e_\lambda)=1$, и нехай
\[
 \Bmod_0=\Amod,\ldots,\Bmod_\lambda=(\Amod,\zeta_\rho,\ldots,\zeta_{\rho-\lambda+1}).
\]
Тогда и $\Bmod_i$ имаје највышши элементарны делитель $e_{\rho-\lambda}$, и зато приходит:
\[
 (e_{\rho-\lambda})=\Bmod_\lambda/\Bmod_{\lambda+1}.
\]}. Формула (6) дале показује, же $\Gmod$ јест делимы чрез квоциент $\Amod/(e_\rho)$; тој квоциент јест делитель $\Amod$ једнакого ранга, зато по дефиницији основного модула од $\Amod$ и обратно јест делимы чрез $\Gmod$. Тымы доказана јест такоже друга формула (5).

Нехай $d$ значи либовольны не идентично изчезајучи детерминант реда $\rho$ из $\Amod$. Тогда $d$ јест делимы чрез $\rho$-ты детерминантны делитель $d_\rho$, и зато чрез $e_\rho$; приходит теды $d\Gmod\zero{\Amod}$, и по горньем закльучаньу следује $\Gmod=\Amod/(d)$. Але за позднејше јест сушчествено, же то последнье квоциентно представјенье за $\Gmod$ може се доказати и без вратаньа к елементарным делительам, и зато имаје обчејшу силу:

\medskip
\noindent\textbf{Теорем III.} \emph{Ако под целыми величинами разумејут се полиномы в неколико неопредельеных}\footnote{В самој ствари користи се само то, же целе величины репродуцирајут се чрез додаванье, одниманье и множенье, при важности обычных правил рачунаньа, значи же они творет ринг; и дале, же в том рингу продукт изчезаје само тогда, когда једен фактор изчезаје, значи же ринг можно разширити до польа творјеньем квоциентов (адјункцијеју паров елементов). Насупрот тому, базисно представјенье модула не користи се; зато (7) важи напр. такоже за област всех целых алгебраичных чисел како коефициентов. (Е. Н.)}, \emph{тогда јеште важи квоциентно представјенье}
\begin{equation}
 \Gmod=\Amod/(d),
\tag{7}
\end{equation}
\emph{кде под $d$ разумеје се либовольны не идентично изчезајучи детерминант реда $\rho$ из $\Amod$.}

Најпрво јест ясно, же појмы ранга, основного модула и модулного квоциента -- кторы тераз уже не става главным идеалом -- при том разширјеньу оставајут зацховане; такоже Теорем I, изкльучивши базисно представјенье.

Нехай ныне
\[
 a_1(\xi)=a_{11}\xi_1+\cdots+a_{1k}\xi_k,\ldots,
 a_\rho(\xi)=a_{\rho1}\xi_1+\cdots+a_{\rho k}\xi_k
\]
будут $\rho$ линеарно независне линеарне формы из $\Amod$, кторе обче не будут творити модулны базис; неопредельене $\xi$ означимо тако, же
\[
 d=|a_{ij}|\ne0;\qquad i,j=1,2,\ldots,\rho.
\]
Из того следује, же $\rho$ линеарных форм специалного облика належет $\Amod$:
\begin{equation}
 d\xi_\lambda+b_{\lambda,\rho+1}\xi_{\rho+1}+\cdots+b_{\lambda,k}\xi_k\zero{\Amod}
 \quad(\lambda=1,2,\ldots,\rho).
\tag{8}
\end{equation}
Але $\Amod$ не може содержати линеарну форму $c_{\rho+1}\xi_{\rho+1}+\cdots+c_k\xi_k$, ктора зависи само од $\xi_{\rho+1},\ldots,\xi_k$, бо иначе најменеј једен детерминант реда $\rho+1$, именно $d\cdot c_\alpha$, был бы различны од нуле.

Ныне квоциент $\Amod/(d)$, како делитель $\Amod$ једнакого ранга, јест делимы чрез $\Gmod$; але важи и обратно. Бо за всако $g(\xi)\zero{\Gmod}$ по дефиницији важи
\[
 c\cdot g(\xi)\zero{\Amod};\quad c\ne0
 \quad\text{и следовательно}\quad d\cdot c\cdot g(\xi)\zero{\Amod}.
\]
По (8) пак
\[
 d\cdot g(\xi)=g_{\rho+1}\xi_{\rho+1}+\cdots+g_k\xi_k\zero{\Amod},
\]
следовательно приходит
\[
 c g_{\rho+1}\xi_{\rho+1}+\cdots+c g_k\xi_k\zero{\Amod}.
\]
По том, что было прáве примечено, следује $c\cdot g_{\rho+1}=0,\ldots,c\cdot g_k=0$, и тым, заради $c\ne0$, такоже $g_{\rho+1}=0,\ldots,g_k=0$; зато $d\cdot g(\xi)\zero{\Amod}$, чим (7) јест доказана.

Треба приметити, же $d\cdot g(\xi)\zero{\Amod}$ следује такоже директно из того, же оба модулы $(a_1(\xi),\ldots,a_\rho(\xi))$ и $(g(\xi),a_1(\xi),\ldots,a_\rho(\xi))$ имајут једнак ранг $\rho$; тогда -- положено $g(\xi)=c_1\xi_1+\cdots+c_k\xi_k$ -- детерминант реда $\rho+1$
\[
\begin{vmatrix}
 a_1(\xi)&a_{11}&\cdots&a_{1\rho}\\
 \vdots&\vdots&&\vdots\\
 a_\rho(\xi)&a_{\rho1}&\cdots&a_{\rho\rho}\\
 g(\xi)&c_1&\cdots&c_\rho
\end{vmatrix}
\]
изчезаје. Але доказ даны выше возврача се в §~4 после формулы (30).


\section*{§ 2. Теоремы о разложеньу норм и модулов.}

Ныне знову јде о модулы линеарных форм взходно к целым рационалным числам или полиномам једној неопредельенеј с коефициентами из $P$.

\medskip
\noindent\textbf{Дефиниция III.} \emph{Ако, како обычно, все линеарне формы в $\xi$, кторе сут конгруентне меджу собоју по модулу $\Amod$, схватити в једну остаткову класу, тогда сымбол $\Gmod\mid\Amod$ означује систем остатковых класов од $\Gmod$ по $\Amod$, то јест систем всех тыцх остатковых класов по $\Amod$, кторе состојут из елементов $\Gmod$. Норма $\Gmod$ взходно к $\Amod$ -- в знаках $\Norm(\Gmod\mid\Amod)$ -- определьа се како детерминант преходного преображеньа од $\Gmod$ к $\Amod$, то јест како детерминант преображеньа, кторо выража либоволны линеарно независны модулны базис $\Amod$ чрез таковы базис основного модула $\Gmod$ од $\Amod$.}

Из (1) и (3), односно из приметкы к Теорему II, добивајемо:
\begin{equation}
 \Norm(\Gmod\mid\Amod)=e_1e_2\cdots e_\rho=d_\rho;
 \qquad
 \Gmod=\Amod/\bigl(\Norm(\Gmod\mid\Amod)\bigr).
\tag{9}
\end{equation}

Из (1), (3) и (9) знаным способом следује, же в случају целых рационалных чисел $\Norm(\Gmod\mid\Amod)$ представјаје число остатковых класов од $\Gmod$ по $\Amod$, дочим в случају полиномов једној неопредельенеј -- кде число остатковых класов обче ставаје бесконечным -- степень $\Norm(\Gmod\mid\Amod)$ ставаје једнак \emph{числу остатковых класов, линеарно независных взходно к $P$}. Ако $\Bmod$ јест делитель $\Amod$ једнакого ранга, тогда $\Bmod$ вместе с представительем либовольној такој остатковеј класы содержи и все елементы тој класы; значи $\Bmod$ настава из $\Amod$ доданьем конечно многих остатковых класов, односно их линеарных комбинациј. Из того одмах следује:

\medskip
\noindent\textbf{Теорем IV.} \emph{Ако $\Bmod$ јест делитель $\Amod$ једнакого ранга, тако же основне модулы совпадајут, и ако наврх того $\Norm(\Gmod\mid\Bmod)=\Norm(\Gmod\mid\Amod)$, тогда такоже $\Bmod=\Amod$.}

О связи меджу разложеньем норм и модулов важи следујуче тврдјенье.

\medskip
\noindent\textbf{Теорем V.} \emph{Нехай}
\begin{equation}
 \Norm(\Gmod\mid\Amod)=r\cdot s,\qquad (r,s)=1;
\tag{10}
\end{equation}
\emph{тогда обстаје једин и само једин модул $\Rmod$, и такоже једин и само једин модул $\Smod$, оба делительи $\Amod$ једнакого ранга, таке, же}
\[
 r=\Norm(\Gmod\mid\Rmod),\qquad s=\Norm(\Gmod\mid\Smod)
\]
\emph{важи. $\Rmod$ и $\Smod$ определьајут се како}
\[
 \Rmod=\Amod/(s);\qquad \Smod=\Amod/(r),
\]
\emph{и $\Amod$ јест једнак најменшему обчему многократному од $\Rmod$ и $\Smod$.}\footnote{Најменше обче многократно $[\Rmod,\Smod]$ ту разумеје се в модулном смыслу: како целост линеарных форм, кторе содержат се и в $\Rmod$ и в $\Smod$. Одповедајуче највечши обчи делитель $(\Rmod,\Smod)$ в модулном смыслу разумеје се како целост линеарных форм, кторе можно представити како суму елемента из $\Rmod$ и елемента из $\Smod$.}
\emph{Ако положимо $r_\rho=(r,e_\rho)$, $s_\rho=(s,e_\rho)$, тогда важе взаимности}
\begin{equation}
 (s_\rho)=\Amod/\Rmod;\quad \Rmod=\Amod/(s_\rho);
 \qquad
 (r_\rho)=\Amod/\Smod;\quad \Smod=\Amod/(r_\rho).
\tag{11}
\end{equation}

Бо ако обчејше положимо $r_i=(r,e_i)$, $s_i=(s,e_i)$, тогда из (9) и (10) приходит:
\[
 (r_i,s_i)=1;\qquad e_i=r_is_i;\qquad r=r_1\cdots r_\rho;\qquad s=s_1\cdots s_\rho,
\]
и всакы $r_i$, односно $s_i$, дели следујучи $r_{i+1}$, односно $s_{i+1}$. Ако теды положимо
\[
 \Rmod=(r_1\eta_1,\ldots,r_\rho\eta_\rho);\qquad
 \Smod=(s_1\eta_1,\ldots,s_\rho\eta_\rho),
\]
тогда $\Rmod$ и $\Smod$ ставајут делительи $\Amod$ једнакого ранга; заради взаимнеј простоты $r_i$ и $s_i$ модул $\Amod$ ставаје једнак најменшему обчему многократному од $\Rmod$ и $\Smod$, и дале приходит
\[
 \Norm(\Gmod\mid\Rmod)=r_1\cdots r_\rho=r;\qquad
 \Norm(\Gmod\mid\Smod)=s_1\cdots s_\rho=s.
\]
Дале важи
\[
 s_\rho\Rmod\zero{\Amod};\qquad r_\rho\Smod\zero{\Amod};
\]
из $c\Rmod\zero{\Amod}$ следује $cr_\rho\eta_\rho\zero{\Amod}$, зато $c\zero{(s_\rho)}$, чим доказано јест $(s_\rho)=\Amod/\Rmod$ и одповедајуче $(r_\rho)=\Amod/\Smod$. Из
\[
 b(\eta)=b_1\eta_1+\cdots+b_\rho\eta_\rho\zero{\Amod/(s_\rho)}
\]
дале следује, заради взаимнеј простоты всех $r_i$ к $s_\rho$, же $b_i\zero{(r_i)}$; значи $\Rmod=\Amod/(s_\rho)$ и одповедајуче $\Smod=\Amod/(r_\rho)$. Тымы доказаны сут взаимности (11), кторе в специалном случају $r=1$ преходет в (5).

Оставаје доказати \emph{једнозначност} $\Rmod$ и $\Smod$. Нехай $\overline{\Rmod}$ и $\overline{\Smod}$ сут модулы, кторе задовольајут условјам Теорема V, а $\bar r_i$ и $\bar s_i$ сут их елементарне делительи. Понеже $\bar r_i$ и $\bar s_i$ сут делительи од $e_i$, из
\[
 r=\bar r_1\cdots\bar r_\rho,
 \qquad
 s=\bar s_1\cdots\bar s_\rho,
 \qquad
 (\bar r_i,\bar s_i)=1
\]
приходит такоже
\[
 e_i=\bar r_i\bar s_i,
 \qquad
 \bar r_i=(r,e_i)=r_i,
 \qquad
 \bar s_i=(s,e_i)=s_i.
\]
Тогда $\overline{\Rmod}$ и $\overline{\Smod}$ совпадајут с $\Rmod$ и $\Smod$ по елементарным делительам и основному модулу. Зато, ако чрез унимодуларно преображенье $\eta=Q(\xi)$ ввести нове неопредельене $\xi$, тогда буде
\[
 \overline{\Rmod}=(r_1\xi_1,\ldots,r_\rho\xi_\rho);
\]
\[
 \Amod=\bigl(r_1s_1(q_{11}\xi_1+\cdots+q_{1\rho}\xi_\rho),\ldots,
 r_\rho s_\rho(q_{\rho1}\xi_1+\cdots+q_{\rho\rho}\xi_\rho)\bigr).
\]
Из $\Amod\zero{\Rmod}$ тогда приходит
\[
 r_i s_i(q_{i1}\xi_1+\cdots+q_{i\rho}\xi_\rho)\zero{\Rmod};
\]
зато $r_is_iq_{ij}\zero{(r_i)}$. Понеже $(r,s)=1$, то даје $r_iq_{ij}\zero{(r_i)}$, или $\overline{\Rmod}\zero{\Rmod}$. Але понеже $\Norm(\Gmod\mid\Rmod)=\Norm(\Gmod\mid\overline{\Rmod})$, по Теорему IV приходит $\Rmod=\overline{\Rmod}$, и одповедајуче $\Smod=\overline{\Smod}$. Тымы Теорем V доказаны јест во всех честјах.


\section*{§ 3. Идеалы из полиномов.}

Нехай $\overline{\mideal}$ буде {\itshape идеал из полиномов} од $n$ неопредельених $y_1,\ldots,y_n$, с коефициентами из либовольного, абстрактно определенного польа $P$; то јест, вместе с $\Phi_1(y)$ и $\Phi_2(y)$ идеал $\overline{\mideal}$ содержи такоже разлико $\Phi_1(y)-\Phi_2(y)$, и вместе с $\Phi(y)$ содержи такоже $A(y)\cdot\Phi(y)$, кде под $A(y)$ разумеје се либовольны полином с коефициентами из $P$\footnote{Појмовно наставајуче названье "идеал" вместо ранеј обче употребльаных названь "модул" или "модул форм" ту уже оказује се нуждно, да бы се разликовало од модулов из линеарных форм. (Е. Н.)}.

Нехай $y$ буде подложено преображеньу с неопредельеными како коефициентами,
\begin{equation}
 y=U(x);\quad \text{например}\quad
 \begin{cases}
 y_1=x_1,\\
 y_2=u_{21}x_1+x_2,\\
 \cdots\\
 y_n=u_{n1}x_1+\cdots+u_{n,n-1}x_{n-1}+x_n,
 \end{cases}
\tag{12}
\end{equation}
и нехай неопредельене $u_{\mu\nu}$ будут присојединьене к области коефициентов полиномов, ктора тым преходит в поле $P(u)$. Чрез то присојединьенье $\overline{\mideal}$ преходит в $\overline{\mideal}_{(u)}$; зато $\overline{\mideal}_{(u)}$ содержи, кроме полиномов $\Phi_\lambda(y)$ из $\overline{\mideal}$, такоже все линеарне комбинације $\sum \alpha_\lambda(u)\Phi_\lambda(y)$ конечно многих $\Phi_\lambda$, кде под $\alpha_\lambda(u)$ разумејут се рационалне функције од $u_{\mu\nu}$ с коефициентами из $P$. Помноженьем с входным полиномом в $u$ можно, без ограниченьа обчности, сматрати $\alpha_\lambda(u)$ како степенове продукты в $u$. Зато важи:
\begin{equation}
 \sum \alpha_\lambda(u)\Phi_\lambda(y)\zero{\overline{\mideal}_{(u)}};
 \qquad
 \Phi_\lambda(y)\zero{\overline{\mideal}}\quad\text{и обратно}.
\tag{13}
\end{equation}
Посредством (12) $\overline{\mideal}_{(u)}$ преходит в идеал $\mideal$ из полиномов в $x_1,\ldots,x_n$ с коефициентами из $P(u)$:
\begin{equation}
 \overline{\mideal}_{(u)}=\mideal\quad \text{посредством}\quad y=U(x).
\tag{14}
\end{equation}
Спојенье релациј (14) и (13) говори следујуче:
\begin{equation}
\begin{aligned}
 &\text{Из } F(x)\zero{\mideal};\\
 &F(x)=\Phi(y)=\sum \alpha_\lambda(u)\Phi_\lambda(y)\\
 &\hspace{2em}=\sum \alpha_\lambda(u)F_\lambda(x)
 \quad\text{следује } F_\lambda(x)\zero{\mideal}.
\end{aligned}
\tag{15}
\end{equation}
дочим из (14) и (15) в обратном смеру знову добива се (13).

\medskip
\noindent\textbf{Дефиниция IV.} {\itshape Идеалы $\mideal$ из полиномов в $x_1,\ldots,x_n$ с коефициентами из $P(u)$, кторе задовольајут релацијам (15), значи кторе посредством $y=U(x)$ наставају из $\overline{\mideal}_{(u)}$, будут назване трансформоваными идеалами.}

Трансформоване идеалы одликујут се ексистованьем регуларных полиномов; {\itshape полином степена $r$ зове се регуларны взходно к $x_i$, ако он содержи член $x_i^r$ с ненулевым коефициентом}. Видно јест, же полиномы $F_i(x)$ сут регуларне взходно к $x_1$. В следујучем буде суштно поймати те трансформоване идеалы како модулы из линеарных форм в некојих степеновых продуктах од $x$. За то треба утврдити појем основного идеала, кторы позднее прејде в основны модул. Најпрво вводимо означенье, кторо од ныне буде всагда држано:

\medskip
\noindent\textbf{Означенье.} {\itshape Под $F^{(i)},f^{(i)},a^{(i)},\ldots$ будут всагда разумены полиномы од $x$, кторе сут свободне од $x_1,\ldots,x_{i-1}$.}

\medskip
\noindent\textbf{Дефиниция V.} {\itshape К всакому идеалу $\mideal$ определьене сут $n$ основне идеалы $\gideal_0,\gideal_1,\ldots,\gideal_{n-1}$ чрез следујучу установу: основны идеал $(i-1)$-ој ступнье $\gideal_{i-1}$ содержи все и само те полиномы $G(x)$, за кторе обстава полином $b^{(i)}$ -- в обче, меньајучи се вместе с $G(x)$ -- так, же}
\begin{equation}
 b^{(i)}G(x)\zero{\mideal};\qquad b^{(i)}\ne0.
\tag{16}
\end{equation}

То, же $\gideal$ в самој ствари сут идеалы, точно одповедаје доказу модулного својства за $\Gmod$ после формулы (4), при чем (16) јест аналог. Идеал $\gideal_i$ јест делимы чрез $\gideal_{i-1}$, понеже всакы полином $b^{(i+1)}$ једночасно јест некой $b^{(i)}$.

За основне идеалы важи:

\medskip
\noindent\textbf{Теорем VI.} {\itshape Основне идеалы трансформованых идеалов сут трансформоване идеалы}\footnote{Теорем VI употребльаје се само в §~7.}.

Доказ опираје се на познаны Dedekind--Мертенсов теорем\footnote{Dedekind: \emph{Über einen arithmetischen Satz von Gauß}, \emph{Mitt. d. deutsch. math. Ges. zu Prag} 1892. Ф. Мертенс: \emph{Über einen algebraischen Satz}, \emph{Бер. d. Ак. d. Виссенсцх. Виен} 101 (1892), стр. 1560--1566. -- Кронецкерово разширјенье Гауссового теорема на полиномы с неопредельеными коефициентами јест непосредньа последица того теорема.}: нехай $a$ и $b$ будут неопредельене, и нехай положено буде
\[
 \sum a_{i_1,\ldots,i_s}z_1^{i_1}\cdots z_s^{i_s}\cdot
 \sum b_{j_1,\ldots,j_s}z_1^{j_1}\cdots z_s^{j_s}
 =
 \sum c_{k_1,\ldots,k_s}z_1^{k_1}\cdots z_s^{k_s},
\]
\[
 \sum i\le \nu,
 \qquad
 \sum j\le \mu,
 \qquad
 \sum k\le \nu+\mu;
\]
дале нехай $\Amod,\Bmod,\mathfrak C$ означујут модулы, одповедајуче изведене из $a,b,c$ помочу целых рационалных чисел. Тогда обстава експонент $q$ так, же важи идентичност:
\begin{equation}
 \Amod^q\Bmod=\Amod^{q-1}\mathfrak C.
\tag{17}
\end{equation}
Тут $\Amod^q$ состоји из степеновых продуктов $q$-ој степени од $a$ и из их целочисленных линеарных комбинациј.

За доказ Теорема VI положимо ныне:
\begin{equation}
\begin{cases}
 G(x)=\Gamma(y)=\sum \alpha_\lambda(u)\Gamma_\lambda(y)=\sum \alpha_\lambda(u)G_\lambda(x),\\
 b^{(i)}(x)=\varphi(y)=\sum \beta_\mu(u)\varphi_\mu(y),
\end{cases}
\tag{18}
\end{equation}
где $\alpha_\lambda(u)$ и $\beta_\mu(u)$ означујут степенове продукты в $u$. Тогда по (16), (14) и (18) важи:
\begin{equation}
 \sum\beta_\mu(u)\varphi_\mu(y)\cdot
 \sum\alpha_\lambda(u)\Gamma_\lambda(y)\zero{\overline{\mideal}_{(u)}};
 \qquad
 \sum\beta_\mu(u)\varphi_\mu(y)\ne0.
\tag{19}
\end{equation}
На левој стране стоји продукт двох полиномов в $u$, кторых коефициенты сут полиномы $\varphi_\mu(y)$ и $\Gamma_\lambda(y)$; коефициенты продукта-полынома в $u$ сут по правој страни (19) полиномы из $\overline{\mideal}$. Ако теды в (17) заменити $a,b,c$ тыми полиномами в $y$, добиваје се
\[
 \left\{\sum\beta_\mu(u)\varphi_\mu(y)\right\}^q\cdot \Gamma_\lambda(y)\zero{\overline{\mideal}_{(u)}},
\]
што по (18) јест равнозначно с:
\begin{equation}
 b^{(i)}(x)^q\cdot G_\lambda(x)\zero{\mideal};
 \qquad
 G_\lambda(x)\zero{\gideal_{i-1}}.
\tag{20}
\end{equation}
Релације (16), (18) и (20) показујут, же {\itshape основне идеалы $\gideal_{i-1}$ в самој ствари сут трансформоване идеалы.}

\section*{22. Обработка К. Хентзелт:\ к теорији полиномных идеалов и резултантов\quad (продолженье)}
\begin{center}
Math. Анн. 88 (1923), с. 53--79
\end{center}

\section*{§ 4. Вез меджу идеалами из полиномов и модулами из линеарных форм.}

Да бы идеал $\mideal$ из полиномов, кторы все дале предполагаје се трансформованым, можно было сматрати како модул $\Mmod_{i-1}$ $(i=1,\ldots,n)$ из линеарных форм, поједине полиномы $F(x)$ треба сматрати како линеарне формы в степеновых продуктах $\xi_\lambda$ од $x_1\ldots x_{i-1}$:
\[
  F(x)=\sum a^{(i)}_\lambda\xi_\lambda,
\]
где, по означеньу из §~3, под $a^{(i)}_\lambda$ разумејут се полиномы в $x_i\ldots x_n$. Из дефиниције идеала одмах следује, же $\Mmod_{i-1}$ има модулну властност взходно к тым целым величинам $a^{(i)},b^{(i)},\ldots$; а понеже бесконечно много степеновых продуктов $\xi$ од $x_1\ldots x_{i-1}$ выступаје само {\itshape линеарно}, они ради својеј линеарној независности играју за $\Mmod_{i-1}$ ролу неопредельеных. Всакы модул $\Mmod_{i-1}$ содержи бесконечно много неопредельеных $\xi$, але в всакој појединој линеарној форме выступаје само конечно много неопредельеных. Такоже всакы основны идеал $\gideal_{i-1}$ треба сматрати како модул линеарных форм $\Gmod_{i-1}$, где неопредельене сут знову степенове продукты од $x_1\ldots x_{i-1}$. За $i=1$ иде само о {\itshape једну} неопредельену $\xi_0$, ктора одповедаје јединици.

Ныне можно, и на том основано јест все следујуче, мимо то бесконечно много неопредельеных $\xi$, ограничити се на модулы линеарных форм в конечно многих неопредельеных, по следујучем теорему.

\medskip
\noindent\textbf{Теорем VII.} {\itshape Систем остатковых класов $\Gmod_{i-1}\mid \Mmod_{i-1}$ јест изоморфны систему остатковых класов $\Gmod^*_{i-1}\mid \Mmod^*_{i-1}$, где $\Mmod^*_{i-1}$ јест модул в конечно многих неопредельеных, а $\Gmod^*_{i-1}$ означује его основны модул. Модул $\Mmod_{i-1}$ после присојединьеньа $x_{i+1}\ldots x_n$ к $P(u)$ има само конечно много елементарных делительев различных од $1$, именно те елементарне делителье $\Mmod^*_{i-1}$, кторе сут различне од $1$.}

При том елементарне делителье $\Mmod_{i-1}$ треба дефинирати како в приметце 8). Изоморфија, ктора вступаје в Теорем VII, основана јест на следујушеј дефиницији.

\medskip
\noindent\textbf{Дефиниция V.} {\itshape Два система остатковых класов зовут се изоморфными, ако можно их взаемно еднозначно одповедати тако, же разлике двох класов одповедаје разлика одповедајучих класов, а продукту класа с целоју величиноју одповедаје продукт одповедајучего класа с тоју самоју целоју величиноју.}

Доказ Теорема VII добива се полноју индукцијеју. За $i=1$ Теорем VII јест очивидны: $\Mmod_0$ јест модул линеарных форм в једнеј неопредельеној $\xi_0$, ктора одповедаје степеновому продукту нултеј степени од $x$, то јест јединици; целе величины сут все полиномы в $x_1\ldots x_n$. Основны идеал $\gideal_0$ јест јединичны идеал, состојаны из всех полиномов в $x_1\ldots x_n$, зато $\Gmod_0$ јест модул $(\xi_0)$, основны модул од $\Mmod_0$, тако же тут $\Gmod^*_0\mid\Mmod^*_0$ совпада с $\Gmod_0\mid\Mmod_0$. Наконец, понеже $\Mmod_0$ има ранг $1$, он може имати највише једен элементарны делитель различны од $1$.

Најпрво прејдемо од $i=1$ к $i=2$, да бы с тут добивеным разуменьем стројеньа $\Gmod_1\mid\Mmod_1$ провести обчи индукцијны крок. К тому најпрво покажемо, же за $\Gmod_0\mid\Mmod_0$ можно взяти систем представительев ограничены по $x_1$; ради делителности $\gideal_1$ чрез $\gideal_0$ то потом приведе к конечно многим неопредельеным модула $\Gmod^*_1$.

По §~3, како трансформованы идеал, $\mideal$ има најменеј једен полином $C^{(1)}(x)$ регуларны взходно к $x_1$, рецимо $k$-ој степени; зато $\Mmod_0$ содержи линеарну форму $C^{(1)}\xi_0$. Ради регуларности $C^{(1)}(x)$ всакы полином $H(x)$ допусча представјенье
\begin{equation}
 H(x)=b^{(2)}_0+b^{(2)}_1x_1+\cdots+b^{(2)}_{k-1}x_1^{k-1}\quad (C^{(1)}(x));
\tag{21}
\end{equation}
зато, взходно к принадлежности $C^{(1)}\xi_0$ к $\Mmod_0$, за $\Gmod_0\mid\Mmod_0$ в самој ствари можно избрати систем представительев, кторы не досега $k$-ој степени по $x_1$.

Ако ныне особливо за полиномы $G(x)$ из $\gideal_1$ и $F(x)$ из $\mideal$ положити:
\begin{equation}
 \left\{\begin{aligned}
  G(x)&=g^{(2)}_0+g^{(2)}_1x_1+\cdots+g^{(2)}_{k-1}x_1^{k-1} &&(C^{(1)}(x)),\\
  F(x)&=a^{(2)}_0+a^{(2)}_1x_1+\cdots+a^{(2)}_{k-1}x_1^{k-1} &&(C^{(1)}(x)),
 \end{aligned}\right.
\tag{22}
\end{equation}
и ако $\Gmod^*_1$ означује систем всех линеарных форм $g(\xi)=g^{(2)}_0\xi_0+\cdots+g^{(2)}_{k-1}\xi_{k-1}$, одповедајуче $\Mmod^*_1$ систем линеарных форм $a(\xi)=a^{(2)}_0\xi_0+\cdots+a^{(2)}_{k-1}\xi_{k-1}$, тогда из идеалној властности $\gideal_1$ и $\mideal$ следује, же $\Gmod^*_1$ и $\Mmod^*_1$ сут модулы из линеарных форм в $\xi_0\ldots\xi_{k-1}$ взходно к целым величинам $b^{(2)},c^{(2)},\ldots$. Такоже главны идеал выведены из $C^{(1)}(x)$ даје модул взходно к тым целым величинам:
\[
 \Gmod_1=(\xi_0,\xi_1,\ldots,\xi_r,\ldots),
\]
где положено $\xi_r=x_1^rC^{(1)}(x)$. Елементы $\xi_r$ сут линеарно независне взходно к тым целым величинам $b^{(2)},c^{(2)},\ldots$, и меджу собоју и од конечно многих $\xi$. Из делителности $\Gmod_1$ чрез $\Mmod_1$ и зато чрез $\Gmod_1$ следује, же такоже $\Gmod^*_1$ јест делимы чрез $\Gmod_1$, односно $\Mmod^*_1$ чрез $\Mmod_1$. С увагојеньем (22) добива се:
\begin{equation}
 \Gmod_1=(\Gmod^*_1,\Gmod_1),\qquad
 \Mmod_1=(\Mmod^*_1,\Gmod_1).
\tag{23}
\end{equation}

Из (23) и линеарној независности $\xi$ и $\xi$ следује Теорем VII за $i=2$ такто. Ради делителности $\Gmod_1$ чрез $\Mmod_1$ најпрво добива се $\Gmod_1\mid\Mmod_1=\Gmod^*_1\mid\Mmod_1$. Да бы доказати изоморфију с $\Gmod^*_1\mid\Mmod^*_1$, нехай некторе линеарне формы $g^*(\xi)$ из $\Gmod^*_1$ дају систем представительев од $\Gmod^*_1\mid\Mmod^*_1$. Ради линеарној независности $\xi$ и $\xi$, из $g^*(\xi)\equiv 0(\Mmod_1)$ следује такоже $g^*(\xi)\equiv 0(\Mmod^*_1)$; формы $g^*(\xi)$ зато сут такоже неконгруентне модуло $\Mmod_1$, творе систем представительев од $\Gmod_1\mid\Mmod_1$, и одповеданье од $\Gmod_1\mid\Mmod_1$ к $\Gmod^*_1\mid\Mmod^*_1$, посредством того система представительев, јест изоморфна по Дефиницији V.

Совсем одповедајуче добива се: из $b^{(2)}g(\xi)\equiv0(\Mmod_1)$, $b^{(2)}\ne0$, следује $b^{(2)}g(\xi)\equiv0(\Mmod^*_1)$, $b^{(2)}\ne0$. Понеже то изполньаје се за все $g(\xi)$ из $\Gmod^*_1$ и само за те -- бо $\Gmod^*_1$ по (23) состоји из всех полиномов основного идеала $\gideal_1$, кторе једночасно сут линеарне формы в $\xi$ -- $\Gmod^*_1$ тым самым карактеризује се како основны модул од $\Mmod^*_1$.

Наконец, присојединьајучи $x_3\ldots x_n$ к $P(u)$, добива се:
\begin{equation}
\left\{\begin{aligned}
 \Gmod^*_1&=(\eta_1,\ldots,\eta_\rho),&
 \Mmod^*_1&=(e_1\eta_1,\ldots,e_\rho\eta_\rho),\\
 \Gmod_1&=(\eta_\rho,\ldots,\eta_1,\xi_0,\ldots,\xi_r,\ldots),&
 \Mmod_1&=(e_\rho\eta_\rho,\ldots,e_1\eta_1,\xi_0,\ldots,\xi_r,\ldots),
\end{aligned}\right.
\tag{24}
\end{equation}
и зато:
\[
 (e_\rho)=\Mmod_1/\Gmod_1=\Mmod_1/(\Mmod_1,\eta_\rho),\qquad
 (e_{\rho-1})=(\Mmod_1,\eta_\rho)/\Gmod_1\quad\hbox{и т. d.}
\]
Тым, с увагојеньем приметкы 8), $e_\rho,e_{\rho-1},\ldots,e_1,1,\ldots,1,\ldots$ познавајут се како елементарне делителье $\Mmod_1$. Зато за $i=2$ Теорем VII доказаны јест во всих честјах.

Треба приметити, же (22) и (23) користајут само то, же $C^{(1)}(x)$ јест регуларны взходно к $x_1$. Те формулы и к ньим привязане размысльеньа, ведуче к Теорему VII, зато оставајут в силе, ако вместо (12) положи се за основу специално преображенье
\begin{equation}
 y_1=x_1;\quad y_2=u_{21}x_1+z_2;\quad\ldots;\quad y_n=u_{n1}x_1+z_n
\tag{25}
\end{equation}
кторо посредством сложеньа с
\begin{equation}
\left\{\begin{array}{l}
 z=U_1(x)\quad\hbox{или}\\[2mm]
 z_2=x_2;\quad z_3=u_{32}x_2+x_3;\quad\ldots;\quad z_n=u_{n2}x_2+\cdots+u_{n,n-1}x_{n-1}+x_n
\end{array}\right.
\tag{26}
\end{equation}
знову даје (12). Ако $\widetilde{\mideal}_{(u)}$ посредством (25) преходи в $\widetilde{\mideal}$, и ако $\widetilde{\gideal}_1,\widetilde{\Gmod}_1,\ldots$ имајут одповедајуче значенье за $\widetilde{\mideal}$ како $\gideal_1,\Gmod_1,\ldots$ за $\mideal$, тогда важи:
\begin{equation}
 \widetilde{\Gmod}_1=(\widetilde{\Gmod}^{*}_1,\widetilde{\Gmod}_1),\qquad
 \widetilde{\Mmod}_1=(\widetilde{\Mmod}^{*}_1,\widetilde{\Gmod}_1).
\tag{27}
\end{equation}

При том (27) настава из (23) просто чрез обраченье (26). То јест очивидно за $\mideal$ и $C^{(1)}(x)$, зато такоже за $\Gmod_1$ и $\Mmod_1$. Но то важи и за $\gideal_1$, бо
\[
 b^{(2)}(x)G(x)\equiv0(\mideal),\quad b^{(2)}\ne0,
 \qquad
 \widetilde b^{(2)}(z)\widetilde G(x,z)\equiv0(\widetilde\mideal),\quad \widetilde b^{(2)}\ne0
\]
меджусобно условльујут једно друго чрез $z=U_1(x)$. Зато $\widetilde\gideal_1=\gideal_1$ посредством (26); зато такоже $\widetilde\Gmod_1=\Gmod_1$, и по дефиницији данеј чрез (22) за $\Gmod_1^*$ и $\Mmod_1^*$ то само важи за $\widetilde\Gmod_1^*$ и $\widetilde\Mmod_1^*$, и зато за цело представјенье (27).

За обчи индукцијны крок нехай важе следујуче предпоставјеньа:
\begin{equation}
\left\{\begin{aligned}
 \Gmod_{i-1}&=(\Gmod_{i-1}^*,\Gmod_{i-1}),&
 \Mmod_{i-1}&=(\Mmod_{i-1}^*,\Gmod_{i-1}),\\
 \Gmod_{i-1}&=(\xi_0,\xi_1,\ldots,\xi_r,\ldots),
\end{aligned}\right.
\tag{28}
\end{equation}
где $\Gmod^*_{i-1}$ и $\Mmod^*_{i-1}$ означујут модулы -- взходно к целым величинам $b^{(i)},c^{(i)},\ldots$ -- из линеарных форм в конечно многих неопредельеных $\xi$, некојих степеновых продуктах од $x_1,\ldots,x_{i-1}$; и где $\xi$ сут линеарно независне и меджу собоју и од $\xi$ взходно к тым целым величинам. Представјенье одповедајуче к (28), и линеарна независност $\xi$ и $\xi$, имајут важит такоже тогда, ако вместо (12) положи се за основу специално преображенье
\begin{equation}
\begin{aligned}
 y_1&=x_1;\quad \ldots;\quad y_{i-1}=u_{i-1,1}x_1+\cdots+x_{i-1};\\
 y_i&=u_{i,1}x_1+\cdots+u_{i,i-1}x_{i-1}+z_i;\quad\ldots;\quad
 y_n=u_{n,1}x_1+\cdots+u_{n,i-1}x_{i-1}+z_n
\end{aligned}
\tag{29}
\end{equation}
кторо може се сложити до (12) посредством
\[
 z=U_{i-1}(x)\quad\hbox{или}\quad
 z_i=x_i;\quad z_{i+1}=u_{i+1,i}x_i+x_{i+1};\quad\ldots;\quad
 z_n=u_{n,i}x_i+\cdots+u_{n,n-1}x_{n-1}+x_n,
\]
и именно то специално представјенье има наставати из (28) просто чрез обраченье $z=U_{i-1}(x)$.

Из предпоставјень (28) и линеарној независности $\xi$ и $\xi$ следује изоморфија $\Gmod_{i-1}\mid\Mmod_{i-1}$ с $\Gmod^*_{i-1}\mid\Mmod^*_{i-1}$, чрез одповеданье к тому самому систему представительев, точно тыми самыми размысльеньами, кторе вышше были привязане к (23). Такоже добива се, же $\Gmod^*_{i-1}$ ставаје основным модулом од $\Mmod^*_{i-1}$ и же $\Mmod_{i-1}$ има само конечно много елементарных делительев различных од $1$, именно те елементарне делителье $\Mmod^*_{i-1}$, кторе сут различне од $1$.

За проведенье индукцијного крока најпрво треба знову показати, же за $\Gmod^*_{i-1}\mid\Mmod^*_{i-1}$, а зато и за $\Gmod_{i-1}\mid\Mmod_{i-1}$, можно взяти систем представительев ограничены по $x_i$. То следује из представјеньа квоциента доказанного в (7), Теорем III:
\begin{equation}
 \Gmod^*_{i-1}=\Mmod^*_{i-1}/(C^{(i)}),
\tag{30}
\end{equation}
где под $C^{(i)}$ разумеје се којиколи не идентично изчезајучи детерминант $\rho$-го степена из $\Mmod^*_{i-1}$, где $\rho$ означује ранг $\Mmod^*_{i-1}$. Из чести предпоставјень, ктора односи се к специалному преображеньу (29), следује, же $C^{(i)}$ можно всегда пријети регуларным взходно к $x_i$. Модул $\widetilde\Mmod^*_{i-1}$ одповедајучи к $\Mmod^*_{i-1}$ има бо еднакы ранг; ако $\widetilde C^{(i)}$ означује којиколи не идентично изчезајучи детерминант $\rho$-го степена из $\widetilde\Mmod^*_{i-1}$, кторы чрез $z=U_{i-1}(x)$ преходи в регуларны $C^{(i)}$, тогда по предпоставјеньу $C^{(i)}$ ставаје детерминантом $\rho$-го степена из $\Mmod^*_{i-1}$.

Из ексистенције $C^{(i)}$ следује, како в (8), при входном нумерованьу $\xi$, принадлежност $\rho$ линеарных форм специалного облика
\[
 L_\lambda(\xi)=C^{(i)}\xi_\lambda+c^{(i)}_{\lambda,1}\xi_{\rho+1}+\cdots+c^{(i)}_{\lambda,s-\rho}\xi_s
 \qquad(\lambda=1\ldots\rho)
\]
к $\Mmod^*_{i-1}$. Ныне можно всаку линеарну форму в $\xi$ приведти к облику:
\begin{equation}
\begin{aligned}
 H(\xi)&=a^{(i)}_1\xi_1+\cdots+a^{(i)}_\rho\xi_\rho\\
 &\quad+b^{(i)}_1\xi_{\rho+1}+\cdots+b^{(i)}_{s-\rho}\xi_s
 \quad (L_1(\xi),\ldots,L_\rho(\xi)),
\end{aligned}
\tag{31}
\end{equation}
где $a^{(i)}_1\ldots a^{(i)}_\rho$ сут ограничене по $x_i$ и не досегајут степена $k$ од $C^{(i)}$. То представјенье јест еднозначно; бо из
\[
 a^{(i)}_1\xi_1+\cdots+a^{(i)}_\rho\xi_\rho+b^{(i)}_1\xi_{\rho+1}+\cdots+b^{(i)}_{s-\rho}\xi_s\equiv0\quad (L_1(\xi),\ldots,L_\rho(\xi))
\]
следује поравнаньем коефициентов при $\xi_1\ldots\xi_\rho$ и с увагојеньем степенов по $x_i$ идентично изчезанье всих $a^{(i)}$ и $b^{(i)}$.

Нехай ныне особливо
\[
 H(\xi)\equiv0(\Gmod^*_{i-1});\qquad\hbox{зато}\qquad C^{(i)}H(\xi)\equiv0(\Mmod^*_{i-1})\quad\hbox{по (30).}
\]
Из того добива се, како при доказу Теорема III,
\[
 C^{(i)}H(\xi)-\sum_{\tau=1}^{s-\rho}\{C^{(i)}b^{(i)}_\tau-\sum_\lambda a^{(i)}_\lambda c^{(i)}_{\lambda,\tau}\}\xi_{\rho+\tau}
 \equiv0(\Mmod^*_{i-1}),
\]
и зато, понеже како при Теорему III коефициенты при $\xi_{\rho+1}\ldots\xi_s$ морајут изчезнути,
\[
 C^{(i)}b^{(i)}_\tau=\sum_\lambda a^{(i)}_\lambda c^{(i)}_{\lambda,\tau}\qquad(\tau=1\ldots s-\rho),
\]
из чего следује, же такоже $b^{(i)}_\tau$ сут ограничене и не досегајут максималного степена од $c^{(i)}_{\lambda,\tau}$. Ексистенција система представительев за $\Gmod^*_{i-1}\mid\Mmod^*_{i-1}$, то јест за $\Gmod_{i-1}\mid\Mmod_{i-1}$, ограниченого по $x_i$ и не досегајучего некојего фиксованого степена $r$, јест тым доказано.

Означимо ныне с $\vartheta_1\ldots\vartheta_t$ конечно многе степенове продукты
\[
 x_i^\alpha\xi_\lambda\quad(\alpha=0,1,\ldots,k-1;\ \lambda=1,2,\ldots,\rho),\qquad
 x_i^\beta\xi_{\rho+\tau}\quad(\beta=0,1,\ldots,r-1;\ \tau=1,2,\ldots,s-\rho)
\]
и упорјадкујмо изчислимо бесконечно много изразов
\[
\begin{aligned}
 &x_i^\gamma L_\lambda\quad(\gamma=0,1,\ldots\text{ до бесконечности};\ \lambda=1,2,\ldots,\rho),\\
 &x_i^\gamma\xi_\nu\quad(\gamma,\nu=0,1,\ldots\text{ до бесконечности}).
\end{aligned}
\]
в просты бесконечны рјад $\omega_0,\omega_1,\ldots,\omega_\mu,\ldots$. Тогда $\vartheta$ и $\omega$ сут линеарно независне, и поједино меджу собоју и једне од других, взходно к целым величинам $b^{(i+1)},c^{(i+1)},\ldots$. Бо из
\[
 \sum c^{(i+1)}_\alpha x_i^\alpha\xi_\lambda+
 \sum c^{(i+1)}_\beta x_i^\beta\xi_{\rho+\tau}+
 \sum d^{(i+1)}_\gamma x_i^\gamma L_\lambda(\xi)+
 \sum e^{(i+1)}_{\gamma\nu}x_i^\gamma\xi_\nu=0,
\]
где всака сума обхвача само конечно много членов, ради предпоставјене линеарној независности $\xi$ и $\xi$, и $\xi$ меджу собоју взходно к целым величинам $b^{(i)}$, следује изчезанье всих $e^{(i+1)}_{\gamma\nu}$. Из једнозначности представјеньа (31) дале следује изчезанье $c^{(i+1)}_\alpha$ и $c^{(i+1)}_\beta$, и из того наконец, ради линеарној независности $L_\lambda(\xi)$, изчезанье $d^{(i+1)}_\gamma$.

Ако ныне модул $(\Gmod_{i-1},L_1(\xi),\ldots,L_\rho(\xi))$ сматрати како модул $\Gmod_i$ взходно к целым величинам $b^{(i+1)}$, тогда $\Gmod_i$ јест делимы чрез $\Mmod_i$, ради делителности $\Gmod_{i-1},L_1(\xi)\ldots L_\rho(\xi)$ чрез $\Mmod_{i-1}$, и добива се
\[
 \Gmod_i=(\omega_0,\omega_1,\ldots,\omega_\mu,\ldots).
\]
За всако $H(x)\equiv0(\gideal_{i-1})$ из (28), (31) и того, што было доказано вышше, важи:
\[
 H(x)=b^{(i+1)}_1\vartheta_1+\cdots+b^{(i+1)}_t\vartheta_t\quad(\Gmod_i)
\]
како модулно равнанье взходно к целым величинам $b^{(i+1)},c^{(i+1)}$. Но $\gideal_i$ јест делимы чрез $\gideal_{i-1}$, а $\mideal$ делимы чрез $\gideal_i$; зато ако особливо за всако $G(x)$ из $\Gmod_i$ и всако $F(x)$ из $\Mmod_i$ положити
\[
 G(x)=g^{(i+1)}_1\vartheta_1+\cdots+g^{(i+1)}_t\vartheta_t\quad(\Gmod_i),\qquad
 F(x)=a^{(i+1)}_1\vartheta_1+\cdots+a^{(i+1)}_t\vartheta_t\quad(\Gmod_i),
\]
и ако модул состојаны из всих $g(\vartheta)$ означити с $\Gmod_i^*$, а модул состојаны из всих $a(\vartheta)$ с $\Mmod_i^*$, тогда точно те саме размысльеньа, кторе ведли к (23), дају:
\begin{equation}
 \Gmod_i=(\Gmod_i^*,\Gmod_i),\qquad
 \Mmod_i=(\Mmod_i^*,\Gmod_i),\qquad
 \Gmod_i=(\omega_0,\omega_1,\ldots,\omega_\mu,\ldots).
\tag{32}
\end{equation}
При изведеньу (32) знову была користана само регуларност $C^{(i)}$ взходно к $x_i$, зато цело размысльенье оставаје в силе, ако специално преображенье (29) положи се за основу за једен индекс выше. Тогда точно те саме размысльеньа, кторе привязујут се к (27), показујут, же то специално представјенье настава из (32) просто чрез обраченье $z=U_i(x)$.

Тым все предпоставјеньа (28) и т. d. обчего индукцијного крока доказане сут за једен индекс выше; а понеже из тыцх предпоставјень, како тут показано, непосредствено следује Теорем VII, Теорем VII доказаны јест обче. Тым самым показано јест, же произвольност пары $(\Gmod^*_{i-1},\Mmod^*_{i-1})$, дана произвольностју избора $C^{(i)}$, знову укида се изоморфијеју система остатковых класов $\Gmod^*_{i-1}\mid\Mmod^*_{i-1}$.

Ако особливо поле $P$, положено в §~3 како област коефициентов, има карактеристику нул -- то јест, ако просто поле, вкльучено в $P$ и изведено из јединице, јест типа рационалных чисел -- тогда неопредельене $u_{\mu\nu}$ можно специализовати в величины $\bar u_{\mu\nu}$ из $P$ тако, же за $i=1\ldots n$ всако $C^{(i)}$ оставаје регуларно взходно к $x_i$. Горне размысльеньа показујут, же Теорем VII оставаје в силе и при тој специализацији.\footnote{Хентзелт вместо либовольного $P$ берет само поле всих комплексных чисел и разматрјаје главно тај последни случай, меджутым он кладе неопредельене в основу само при правеј елиминационној теорији. Хентзелт потом дале показује, и то суштно с закльучками тут даными, же за творјенье резултантној формы достаточно јест в всаком $\Mmod_{i-1}$ идти само до некојего конечного степена в $x$, кторы превышаје фиксовану границу. Но недостаје изоморфија $\Gmod_{i-1}\mid\Mmod_{i-1}$ с $\Gmod^*_{i-1}\mid\Mmod^*_{i-1}$ дана в Теорему VII, ктора објашньује независност од степеновых чисел. (Е. Н.)} Основны модул и елементарне делителье при том не морајут нужно наставати из обчехо случаја чрез специализацију $u_{\mu\nu}=\bar u_{\mu\nu}$.\footnote{За $\mideal=(x^2,ux+y)$ добива се $\gideal_0=(1)$, $\gideal_1=(1)$. Ако избрати $C^{(1)}=x^2$, тогда $\Gmod_1^*=(1,x)$, $\Mmod_1^*=(ux+y,xy)$, при чем $\Mmod_1^*$ има елементарне делителье $y^2$ и $1$, и можно избрати $C^{(2)}=y^2$. За $u=0$ $C^{(1)}$ и $C^{(2)}$ оставајут регуларне взходно к $x$ односно $y$; але $\Mmod_1^*=(y,xy)$ има елементарне делителье $y,y$. Зато $\Gmod_1\mid\Mmod_1$ ставаје такоже изоморфны систему остатковых класов конечного модула, але не изоморфны к $\Gmod_1\mid\Mmod_1$. Ако бы се напротив избрало $C^{(1)}=ux+y$ и специализовало тако, же $C^{(1)}$ остало регуларно, тогда бы такоже изоморфија была сохраньена. (Е. Н.)}

\clearpage
\section*{§ 5. Резултантна форма и форма елементарных делительев идеала из полиномов.}

Нехай $x_{i+1}\ldots x_n$ будут присојединьене к $P(u)$, тако же $\Gmod_{i-1}$ и $\Mmod_{i-1}$, а следовательно такоже $\Gmod^*_{i-1}$ и $\Mmod^*_{i-1}$, преходет в модулы линеарных форм, где целе величины можно сматрати како полиномы в {\itshape једној} неопределенеј $x_i$. По Теорему VII в том случају елементарне делителье $\Mmod_{i-1}$, различне од $1$, и највише конечно много елементарных делительев $1$ модула $\Mmod_{i-1}$ совпадајут с теми од $\Mmod^*_{i-1}$. Продукт тыцх елементарных делительев, $\rho$-ты детерминантны делитель $\Mmod^*_{i-1}$, был једнак детерминанту преходного преображеньа од $\Gmod^*_{i-1}$ к $\Mmod^*_{i-1}$, зато по Дефиницији III једнак $N(\Gmod^*_{i-1}\mid\Mmod^*_{i-1})$. По (24), односно по аналогној формули за обче $i$, ктора следује из (28), показује се, же продукт тыцх елементарных делительев такоже ставаје једнак детерминанту преходного преображеньа од $\Gmod_{i-1}$ к $\Mmod_{i-1}$; зато его треба означити с $N(\Gmod_{i-1}\mid\Mmod_{i-1})$. Зато добива се
\[
 N(\Gmod_{i-1}\mid\Mmod_{i-1})=N(\Gmod^*_{i-1}\mid\Mmod^*_{i-1})=R^{(i)}(x),
\]
где полином $R^{(i)}(x)$, то јест највечши обчи делитель в полиномном смыслу всих $\rho$-рядных детерминантов из $\Mmod^*_{i-1}$, можно сматрати како целы и примитивны по $x_{i+1}\ldots x_n$ и зато, како делитель $C^{(i)}$, регуларны взходно к $x_i$. Такоже највышши элементарны делитель $E^{(i)}(x)$ модула $\Mmod_{i-1}$, односно $\Mmod^*_{i-1}$, можно сматрати како целы и примитивны по $x_{i+1}\ldots x_n$ и зато регуларны взходно к $x_i$. То води к следујушеј дефиницији.

\medskip
\noindent\textbf{Дефиниция VI.} {\itshape Полином $R^{(i)}(x)$ означује се како резултант $i$-го нивоја идеала $\mideal$, а $E^{(i)}(x)$ како элементарны делитель $i$-го нивоја. Резултант $i$-го нивоја јест зато норма основного идеала $\gideal_{i-1}$ взходно к $\mideal$, сматрана како норма модула $\Gmod_{i-1}$ взходно к $\Mmod_{i-1}$, при чем $x_{i+1}\ldots x_n$ сут присојединьене к $P(u)$; такоже $E^{(i)}(x)$ при том самом присојединьеньу ставаје једнак квоциенту $\Mmod_{i-1}/\Gmod_{i-1}$. Како резултантна форма односно форма елементарных делительев од $\mideal$ означујут се продукты}
\[
 R_\mideal=R^{(1)}R^{(2)}\cdots R^{(n)},\qquad
 E_\mideal=E^{(1)}E^{(2)}\cdots E^{(n)}.
\]

За резултантну форму и форму елементарных делительев важи следујуче.

\medskip
\noindent\textbf{Теорем VIII.} {\itshape Резултантна форма јест делима чрез форму елементарных делительев; обратно, некоја потяга од $E_\mideal$ јест делима чрез $R_\mideal$. $E_\mideal$ и $R_\mideal$ сут делимы чрез $\mideal$; обчеје, $E^{(i)}\cdots E^{(n)}\gideal_{i-1}\equiv0(\mideal)$ и зато $R^{(i)}\cdots R^{(n)}\gideal_{i-1}\equiv0(\mideal)$.}

Именно, понеже норма јест делима чрез највышши элементарны делитель и обратно некоја потяга того највышшего елементарного делительа јест делима чрез норму, та делителност следује за $R^{(i)}(x)$ и $E^{(i)}(x)$ при присојединьеньу $x_{i+1}\ldots x_n$. Але $R^{(i)}(x)$ и $E^{(i)}(x)$ предполагајут се примитивными полиномами взходно к $x_{i+1}\ldots x_n$, зато делителност важи взходно к всим неопредельеным $x_i,x_{i+1},\ldots,x_n$. Так по дефиницији тыцх изразов взходно к всим неопредельеным $x$ важи делителност $R_\mideal$ чрез $E_\mideal$ и некојеј потяги $E_\mideal$ чрез $R_\mideal$.

Дале, по Теорему VII и Теорему II, при присојединьеньу $x_{i+1}\ldots x_n$, највышши элементарны делитель $E^{(i)}(x)$ дефинованы јест како квоциент $\Mmod_{i-1}/\Gmod_{i-1}$. Зато, ако знову прејдемо к полиномам в всих неопредельеных $x$, за всако
\begin{equation}
 G(x)\equiv0(\gideal_{i-1})
 \quad\hbox{јест}\quad b^{(i+1)}\ne0,
 \quad\hbox{тако же}\quad b^{(i+1)}E^{(i)}(x)G(x)\equiv0(\mideal).
\tag{33}
\end{equation}
Ныне особливо $\gideal_0$ ставаје јединичным идеалом, зато добива се:
\[
 b^{(2)}E^{(1)}(x)\equiv0(\mideal),\quad b^{(2)}\ne0,
 \qquad\hbox{и зато}\qquad E^{(1)}(x)\equiv0(\gideal_1).
\]
Обче, по (33), применьеној к
\[
 E^{(1)}(x)\cdots E^{(i-1)}(x)\equiv0(\gideal_{i-1}),
\]
јест $b^{(i+1)}\ne0$ тако, же $b^{(i+1)}E^{(1)}\cdots E^{(i)}\equiv0(\mideal)$; зато $E^{(1)}\cdots E^{(i)}\equiv0(\gideal_i)$. Понеже $\gideal_n=\mideal$, добива се
\begin{equation}
 E_\mideal=E^{(1)}\cdots E^{(n)}\equiv0(\mideal)
 \quad\hbox{и зато}\quad
 R_\mideal=R^{(1)}\cdots R^{(n)}\equiv0(\mideal).
\tag{34}
\end{equation}

Ако $G(x)$ в (33) преходи чрез конечно много полиномов некојеј базы идеала $\gideal_{i-1}$ и $c^{(i+1)}(x)$ означује продукт одповедајучих $b^{(i+1)}(x)$, тогда добива се
\[
 c^{(i+1)}E^{(i)}(x)\gideal_{i-1}\equiv0(\mideal);
 \qquad c^{(i+1)}\ne0
 \quad\hbox{и зато}\quad E^{(i)}(x)\gideal_{i-1}\equiv0(\gideal_i),
\]
и из того, како выше, чрез конечно повторјенье:
\[
 E^{(n)}(x)\cdots E^{(i)}(x)\gideal_{i-1}\equiv0(\mideal),
\]
и зато такоже
\[
 R^{(n)}(x)\cdots R^{(i)}(x)\gideal_{i-1}\equiv0(\mideal),
\]
чим Теорем VIII јест доказаны во всих честјах.

Теорем VIII основан јест суштно на властностјах формы елементарных делительев; але же резултантна форма има карактеристично значенье за $\mideal$, показује

\medskip
\noindent\textbf{Теорем IX.} {\itshape Ако $\nideal$ јест делитель од $\mideal$ -- делитель разумены в идеалном смыслу -- и ако резултантне формы $R_\nideal$ и $R_\mideal$ совпадајут, тогда идеалы $\nideal$ и $\mideal$ совпадајут.}

Именно, нехай $\gideal_0\ldots\gideal_{n-1},\gideal_n=\mideal$ и $\hideal_0\ldots\hideal_{n-1},\hideal_n=\nideal$ означујут основне идеалы од $\mideal$ и $\nideal$. Најпрво $\gideal_0$ и $\hideal_0$ ставајут оба јединичным идеалом, зато $\gideal_0=\hideal_0$. Ако ныне предпоставити $\gideal_{i-1}=\hideal_{i-1}$, тако же $\Mmod_{i-1}$ и $\Nmod_{i-1}$ имајут тој самы основны модул $\Gmod_{i-1}$, тогда предпоставјенье $R_\nideal=R_\mideal$ при присојединьеньу $x_{i+1}\ldots x_n$ можно записати како:
\[
 N(\Gmod_{i-1}\mid\Nmod_{i-1})=N(\Gmod_{i-1}\mid\Mmod_{i-1})
 \quad\hbox{или такоже}\quad
 N(\Gmod^*_{i-1}\mid\Nmod^*_{i-1})=N(\Gmod^*_{i-1}\mid\Mmod^*_{i-1}),
\]
при том положено јест одповедајуче к (32): $\Nmod_{i-1}=(\Nmod^*_{i-1},\Gmod_{i-1})$, тако же $\Nmod^*_{i-1}$ такоже ставаје модулом линеарных форм в $\xi$, а $\Gmod^*_{i-1}$ его основным модулом. Ради делителности $\Mmod_{i-1}$ чрез $\Nmod_{i-1}$, ктора води к делителности $\Mmod^*_{i-1}$ чрез $\Nmod^*_{i-1}$, из того по Теорему IV следује, же при том присојединьеньу два модулы $\Mmod^*_{i-1}$ и $\Nmod^*_{i-1}$, и зато такоже $\Nmod_{i-1}$ и $\Mmod_{i-1}$, совпадајут. Присојединьенье $x_{i+1}\ldots x_n$ значи але, же к $\Mmod_{i-1}$ односно к $\mideal$ додавајут се все такие полиномы $G(x)$, за кторе
\[
 b^{(i+1)}(x)G(x)\equiv0(\mideal),
 \qquad b^{(i+1)}\ne0
\]
важи, то јест, чрез присојединьенье $\Mmod_{i-1}$ односно $\mideal$ преходи в $\gideal_i$, одповедајуче $\nideal$ в $\hideal_i$; зато $\gideal_i=\hideal_i$ и следовательно $\gideal_n=\hideal_n$, то јест $\mideal=\nideal$ доказано.

Наконец то само начело преноса, применьено к Теорему V, даје још

\medskip
\noindent\textbf{Теорем X.} {\itshape Нехай $R^{(i)}=S^{(i)}T^{(i)}$ буде разложенье $R^{(i)}$ на -- в полиномном смыслу -- взаимно просте полиномы $S^{(i)}$ и $T^{(i)}$. Тогда ексистује једен и само једен идеал $\jideal$, такоже једен и само једен идеал $\tideal$, оба делителье од $\mideal$ с једнакым основным идеалом $(i-1)$-го нивоја $\gideal_{i-1}$, тако же $S^{(i)}$ ставаје једнак $N(\Gmod_{i-1}\mid\Smod_{i-1})$, а $T^{(i)}$ једнак $N(\Gmod_{i-1}\mid\mathfrak T_{i-1})$. $\jideal$ и $\tideal$ дефиноване сут како целост полиномов $S(x)$ односно $T(x)$ такых, же}
\[
 s^{(i+1)}(x)T^{(i)}(x)S(x)\equiv0(\mideal),\quad s^{(i+1)}\ne0,
\]
\[
 t^{(i+1)}(x)S^{(i)}(x)T(x)\equiv0(\mideal),\quad t^{(i+1)}\ne0,
\]
{\itshape и $\gideal_i$ ставаје најменше обче многократно од $\jideal$ и $\tideal$,}
\[
 \gideal_i=[\jideal,\tideal].
\]

К тому треба само приметити, же $s^{(i+1)},t^{(i+1)}$ за $\jideal$ и $\tideal$ можно избрати фиксовано, како продукт $s^{(i+1)},t^{(i+1)}$, кторе одповедајут базным елементам од $\jideal$ односно $\tideal$, и же, како выше примечено, при присојединьеньу $x_{i+1}\ldots x_n$ $\mideal$ преображаје се в $\gideal_i$. Особливо разложенье $R^{(i)}$ можно продолжити до потяг неразложимых полиномов -- примарных факторов -- што води к одповедајучему представјеньу $\gideal_i$ како најменшего обчего многократного.

\clearpage
\section*{§ 6. Властна теория елиминације.}

По Теорему VIII односно по формули (34) резултантна форма и форма елементарных делительев сут делимы чрез $\mideal$, зато изчезајут за все нулы $\mideal$; при том под ``нулами $\mideal$'' разумејут се таке -- приналежне алгебраично затворјеному польу, выведеному из $P(u;x_{i+1}\ldots x_n)$\footnote{Же всако поле, и то в суштно једнозначном способу, можно разширити до алгебраично затворјеного, показал \mbox{Steinitz} в својој \emph{Algebraische Theorie der Körper} (\mbox{J. f. M.} 137 (1910), стр. 167--309), и тым дал рационалны еквивалент фундаменталному теорему алгебры. В самој ствари за всако $i=1,\ldots,n$ јде о конечно разширительно поле од $P(u,x_{i+1},\ldots,x_n)$. (Е. Н.)} -- системи величин $x_1\ldots x_i$, же за те системи величин всакы полином из $\mideal$ изчезаје, при чем $x_{i+1}\ldots x_n$ мысльет се присојединьене к области коефициентов $(i=1,\ldots,n)$. Те присојединьене $x_{i+1}\ldots x_n$, како буде показано, можно такоже заменити величинами из $P(u)$. Содержанье теорије елиминације, обратно, јест выводженье нул $\mideal$ из нул резултантној формы. То обратно выводженье основано јест на сукцесивној елиминацији, ктора буде развијена в том параграфу; когда в следујучем параграфу буде доказано, же при том наставајуча форма $D^{(i)}(x)$ јест делима чрез $E^{(i)}(x)$, из делительности некојеј потяги $E^{(i)}(x)$ чрез $R^{(i)}(x)$ следује желано обратно тврдјенье.

Теория сукцесивној елиминације, ктора тут ма быти дана, основана јест на следујучем.

\medskip
\noindent\textbf{Теорем XI.} {\itshape Нехай $\bideal$ означује идеал из полиномов в једнеј неопредельеној $t$ с коефициентами из $P$, значи главны идеал; нехай $h(t)$ буде полином из $\bideal$ степена $k$, а $\Bmod$ модул линеарных форм в $1,t,\ldots,t^{k-1}$, в кторы $\bideal$ преходит модуло $h(t)$. Тогда $\Bmod$ има ранг $k-p$, ако $p$ означује степень базного полинома $f(t)$ идеала $\bideal$.}

По предпоставјеньу $h(t)=h_1(t)f(t)$, где $h_1(t)$ јест степена $k-p$. Зато остаткове класы идеала $\bideal$ модуло $h(t)$, представјене полиномами $f(t),tf(t),\ldots,t^{k-p-1}f(t)$, сут линеарно независне; и всака остаткова класа идеала $\bideal$ модуло $h(t)$ можно линеарно, с коефициентами из $P$, представити чрез те $k-p$ специалне класы. Але $\bideal$ модуло $h(t)$ преходит в модул линеарных форм в $1,t,\ldots,t^{k-1}$, кторого ранг тым самым јест точно $k-p$.

Нехай ныне $\aideal=\aideal_1$ буде трансформованы идеал из полиномов, $A^{(1)}(x)$ полином из $\aideal_1$, регуларны взходно к $x_1$, степена $k_1$, а $\mathfrak A_1$ модул линеарных форм в $1,x_1,\ldots,x_1^{k_1-1}$, с полиномами $a^{(2)}(x)$ како коефициентами, в кторы $\aideal_1$ преходит модуло $A^{(1)}(x)$. Дале нехай $\aideal_2$ означује идеал в $x_2\ldots x_n$, выведены из всих $k_1$-рядных детерминантов из $\mathfrak A_1$. По Теорему XI $\aideal_2$ тогда и само тогда јест једнак нуловому идеалу, когда полиномы из $\aideal_1$ имајут обчи делитель -- делитель разумены в полиномном смыслу -- степена $p_1>0$ взходно к $x_1$. Ако $\aideal_2$ јест различны од нулового идеала, тогда, понеже $\aideal$ был предпоставјен како трансформованы идеал, он содержи полином $A^{(2)}(x)$ степена $k_2$, регуларны взходно к $x_2$, како то непосредно видно из сложеньа трансформације (12) из специалных трансформациј (25) и (26). Нехай знову $\mathfrak A_2$ означује модул линеарных форм в $1,x_2,\ldots,x_2^{k_2-1}$, в кторы $\aideal_2$ преходит модуло $A^{(2)}(x)$, а $\aideal_3$ идеал в $x_3\ldots x_n$, выведены из всих $k_2$-рядных детерминантов из $\mathfrak A_2$, кторы муси содержати полином $A^{(3)}(x)$ регуларны взходно к $x_3$.

Такым способом поступак продолжује се, доколе не приде се к нуловому идеалу или доколе $\aideal_{n+1}$ не стане јединичным идеалом; бо, понеже $\aideal_{n+1}$ јест свободны од $x$, он може быти само нуловы или јединичны идеал. Показује се, же идеалы $\aideal_1,\aideal_2,\ldots$ сут еднозначно определьене чрез $\aideal$, независно од полиномов $A^{(\lambda)}(x)$ положених в основу. То ясно јест за $\aideal_1=\aideal$, зато можно то предпоставити за $\aideal_1,\ldots,\aideal_i$. Ако ныне $\aideal_i$ модуло $A^{(i)}(x)$ преходит в модул линеарных форм $\mathfrak A_i$, тогда $\mathfrak A_i$ можно такоже карактеризовати како целост полиномов из $\aideal_i$, кторе в $x_i$ не достигајут степена $k_i$; зато $\mathfrak A_i$ зависи само од степена $k_i$ полинома $A^{(i)}(x)$. Нехай ныне $\bar A^{(i)}(x)$ буде полином из $\aideal_i$, регуларны взходно к $x_i$, степена $\bar k_i>k_i$, $\overline{\mathfrak A}_i$ одповедајучи модул линеарных форм, а $\bar\aideal_{i+1}$ идеал выведены из $\bar k_i$-рядных детерминантов из $\overline{\mathfrak A}_i$. Тогда важи модулно равнанье
\[
 \overline{\mathfrak A}_i=(\mathfrak A_i,A^{(i)},x_iA^{(i)},\ldots,x_i^{\bar k_i-k_i-1}A^{(i)}).
\]
Понеже, ради регуларности $A^{(i)}(x)$ взходно к $x_i$, полиномы $A^{(i)},x_iA^{(i)},\ldots$ можно ввести како нове неопредельене линеарных форм, из того следује совпаданье $k_i$-рядных детерминантов из $\mathfrak A_i$ с $\bar k_i$-рядными из $\overline{\mathfrak A}_i$; и зато $\bar\aideal_{i+1}=\aideal_{i+1}$.\footnote{То јест в начелу тој самы закльучок, на ктором основана была изоморфија $\Gmod^*_{i-1}\mid\Mmod^*_{i-1}$ с $\Gmod_{i-1}\mid\Mmod_{i-1}$.}

Дале $\aideal_{i+1}$ всагда јест делимы чрез $\aideal_i$. То ясно јест, ако $\aideal_{i+1}$ јест нуловы идеал; в противном случају $\mathfrak A_i$ јест ранга $k_i$, зато $\aideal_{i+1}$ јест по дефиницији идеал, состојаны из $k_i$-рядных детерминантов из $\mathfrak A_i$, делимы чрез $\mathfrak A_i$ и следовательно чрез $\aideal_i$. Всакы $\aideal_{i+1}$ тым самым јест делимы чрез $\aideal$; ако дакле $\aideal_{n+1}$ ставаје јединичны идеал, тогда и $\aideal$ јест јединичны идеал.

Нехай ныне $\aideal$ буде различны од јединичного идеала; $\aideal_1\ne0,\ldots,\aideal_i\ne0;\ \aideal_{i+1}=0$. Нехай $D^{(i)}(x)$ буде највечши обчи делитель -- делитель разумены в полиномном смыслу -- всих полиномов из $\aideal_i$, ексистујучи по Теорему XI; како делитель од $A^{(i)}(x)$ сам $D^{(i)}(x)$ јест регуларны взходно к $x_i$ степена $p_i>0$. Присојединимо $x_{i+1}\ldots x_n$ к $P(u)$; тогда $D^{(i)}(x)$ в алгебраичном разширительнем польу од $P(u;x_{i+1}\ldots x_n)$ допусча разложенье на $p_i$ линеарных факторов. Нехай $x_i-\bar x_i$ буде једен таковы линеарны фактор, тако же за $x_i=\bar x_i$ все полиномы из $\aideal_i$ изчезајут и следовательно по Теорему XI полиномы из $\aideal_{i-1}$ при тој специализацији добивајут највечши обчи делитель $D^{(i-1)}(x)$. Како делитель од $A^{(i-1)}(x)$ сам $D^{(i-1)}(x)$ знову јест регуларны взходно к $x_{i-1}$ степена $p_{i-1}>0$; значи он особливо не може идентично изчезнути и в разширительнем польу од $P(u,\bar x_i,x_{i+1},\ldots,x_n)$ допусча разложенье на $p_{i-1}$ линеарных факторов. Ако $x_{i-1}-\bar x_{i-1}$ јест таковы линеарны фактор, ему одповедаје највечши обчи делитель из $\aideal_{i-2}$. Продолжујучи тако, приходимо к конечно многим обчим нульам $\aideal$, кторе лежат в алгебраичном разширительнем польу од $P(u;x_{i+1}\ldots x_n)$. Обратно, ради делительности $\aideal_i$ чрез $\aideal$, всака нула $\aideal$ јест такоже нула $\aideal_i$. Ако особливо $x_1=\bar x_1,\ldots,x_i=\bar x_i,x_{i+1}=x_{i+1},\ldots,x_n=x_n$ јест нула $\aideal$, из того следује $\aideal_{i+1}=0$; и полиномы из $\aideal_i$ добивајут највечши обчи делитель с линеарным фактором $x_i-\bar x_i$. В суме добива се:

\medskip
\noindent\textbf{Теорем XII.} {\itshape Нехай $\aideal$ буде различны од јединичного идеала, $\aideal_1\ne0,\ldots,\aideal_i\ne0;\ \aideal_{i+1}=0$; тогда при присојединьеньу $x_{i+1}\ldots x_n$ к $P(u)$ ексистује конечно много приналежных системов величин $\bar x_1\ldots\bar x_i$, кторе лежат в алгебраичном разширительнем польу од $P(u,x_{i+1},\ldots,x_n)$, так же все полиномы из $\aideal$ изчезајут за $x_1=\bar x_1,\ldots,x_i=\bar x_i,x_{i+1}=x_{i+1},\ldots,x_n=x_n$. Обратно, ако дана јест така нула $\aideal$, из того следује $\aideal_{i+1}=0$ и наставанье обчего делительа $D^{(i)}(x)$ всих полиномов из $\aideal_i$, кторы има линеарны фактор $x_i-\bar x_i$.}

Теорем XII показује особливо, же всакы идеал $\aideal$, кторы не има нулы, ставаје јединичны идеал.

Да бы провести разложенье, кторо такоже експлицитно показује приналежност системов величин, вводет се нове неопредельене $v$, кторе можно присојединити к $P(u,x_{i+1},\ldots,x_n)$. Положимо:
\begin{equation}
 z_i=-v_1x_1-\cdots-v_{i-1}x_{i-1}+x_i,
\tag{35}
\end{equation}
тако же сложенье (12) с (35) знову дава субституцију типа (12), где ныне само $u_{ik}$ заменьене сут чрез $w_{ik}=u_{ik}-v_k$, а обчеје $u_{i+h,k}$ чрез $w_{i+h,k}=u_{i+h,k}-u_{i+h,i}v_k$. Ако так субституција (35) преводит по предпоставјеньу трансформованы идеал $\aideal$ в $\bideal$, тогда $\bideal$ настава просто из $\aideal$ заменоју $u_{\mu\nu}$ чрез $w_{\mu\nu}$ и присојединьеньем $v$; и тым самым процесом идеалы $\bideal_1,\bideal_2,\ldots$, дефиноване посредством $\bideal$, наставајут из $\aideal_1,\aideal_2,\ldots$. Обратно, $\aideal_i$ настава из $\bideal_i$ посредством $v=0$; идеалы $\aideal_i$ и $\bideal_i$ зато сут једночасно нулове идеалы или једночасно различне од нулового идеала.

Нехай ныне знову $\aideal_1\ne0;\ldots;\aideal_i\ne0;\aideal_{i+1}=0$; зато такоже $\bideal_1\ne0;\ldots;\bideal_i\ne0;\bideal_{i+1}=0$; нехай $\bar x_1\ldots\bar x_i,x_{i+1}\ldots x_n$ буде приналежны систем нул $\aideal$. Ако дефинујемо $\bar z_i=\bar x_i+v_1\bar x_1+\cdots+v_{i-1}\bar x_{i-1}$, тогда $\bar x_1\ldots\bar x_{i-1},\bar z_i,x_{i+1}\ldots x_n$ посредством (35) ставаје нула $\bideal$. Ако дакле $H^{(i)}(z,x)$ означује највечши обчи делитель всих полиномов из $\bideal_i$, кторы можно предполагати целым и примитивным по $v$, тогда по последньеј чести Теорема XII $H^{(i)}(z,x)$ добива фактор
\[
 z_i-\bar z_i=z_i-(\bar x_i+v_1\bar x_1+\cdots+v_{i-1}\bar x_{i-1}),
\]
и всакој нули $\aideal$, ктора настава при присојединьеньу $x_{i+1}\ldots x_n$, одповедаје таковы фактор. Треба показати, же тым разложенье $H^{(i)}(z,x)$ јест изчрпано; то јест, же из $H^{(i)}$ не можно одделити фактор $H^{(i)}_1$, кторы не разклада се на линеарне факторы в $z$ \emph{и} $v$. Нехай то буде случай. Тогда, ради предпоставјенеј примитивности по $v$, $H^{(i)}$ содержи најменеј једен линеарны фактор $z_i-\bar z_i$, где $\bar z_i$ приналежи алгебраичному разширительнему польу од $P(u,v,x_{i+1},\ldots,x_n)$. Ако $\bar x_1\ldots\bar x_{i-1},\bar z_i,x_{i+1}\ldots x_n$ јест одповедајуча нула $\bideal$, она муси совпадати с једној из конечно многих нул $\aideal$, кторе наставајут при присојединьеньу $x_{i+1}\ldots x_n$, и зато муси быти независна од $v$. Тым самым $\bar z_i$ нужно содержи $v$ линеарно, не алгебраично или рационално-нелинеарно; проти предпоставјеньу, из $H^{(i)}_1$ можно одделити јеште једен линеарны фактор $z_i-(\bar x_i+v_1\bar x_1+\cdots+v_{i-1}\bar x_{i-1})$ в $z$ и $v$. Зато как експлицитно разложенье добива се:
\[
 H^{(i)}(z,x)=\prod \{z_i-(\bar x_i+v_1\bar x_1+\cdots+v_{i-1}\bar x_{i-1})\}^{\alpha_\lambda}.
\]
Понеже степень $H^{(i)}$ и поједине експоненты $\alpha_\lambda$ мусет совпадати с одповедајучими од $D^{(i)}(x)$ -- бо $H^{(i)}$ настава из $D^{(i)}$ специализацијеју $u_{\mu\nu}=w_{\mu\nu}$, а $D^{(i)}$ из $H^{(i)}$ специализацијеју $v=0$ -- то разложенье јеште показује, же ради присојединьеньа неопредельеных $u_{\mu\nu}$ к $P$ такоже при елиминацији, ктора выходит из $\aideal$, всакој нули $\bar x_i$ полинома $D^{(i)}$ одповедаје одповедајучи линеарны највечши обчи делитель из $\aideal_{i-1}\ldots\aideal_1$ или потяга такового; зато всакы раз добива се само једен приналежны систем нул $\bar x_1\ldots\bar x_i,x_{i+1}\ldots x_n$ идеала $\aideal$.\footnote{Треба указати, же тут дана сукцесивна елиминација -- в противности к Кронецкеровој методе -- зависи само од даного идеала $\aideal$ и јест независна од всакој идеалној базы; особливо то важи такоже за експоненты $\alpha_\lambda$. Полно проведенье елиминације по тој методе, по Хентзелту, требовало бы продолженье поступка на квоциент $\aideal/D^{(i)}$, што можно опустити с обзиром на следујучи параграф. (Е. Н.)}

\clearpage
\section*{§ 7. Резултантна форма и теория елиминације.}

Да бы установити вез меджу резултантној формоју и теоријеју елиминације, применимо сукцесивну елиминацију из §~6 специално к квоциенту $\aideal=\mideal/\gideal_{i-1}$ идеала чрез основны идеал $(i-1)$-го нивоја. Најпрво треба показати, же тој квоциент представльа трансформованы идеал. Ныне $\mideal$ јест трансформованы, и следовательно по Теорему VI из §~3 такоже $\gideal_{i-1}$. Из
\[
 H(x)\equiv0(\mideal/\gideal_{i-1});\qquad
 H(x)=\overline H(y)=\sum\alpha_i(u)\overline H_i(y)=\sum\alpha_i(u)H_i(x)
\]
следује дакле:
\[
 \overline H(y)\overline\gideal_{i-1,(u)}\equiv0(\overline\mideal_{(u)});
 \quad\hbox{зато такоже}\quad
 \overline H(y)\overline\gideal_{i-1}\equiv0(\overline\mideal_{(u)}),
\]
или
\[
 \overline H_i(y)\overline\gideal_{i-1}\equiv0(\overline\mideal),
 \quad\hbox{зато}\quad H_i(x)\gideal_{i-1}\equiv0(\mideal),
\]
чим доказано јест делительност $H_i(x)$ чрез $\mideal/\gideal_{i-1}$ и тым самым тој идеал како трансформованы, тако же метода елиминације из §~6 ставаје применима.

Але (30), с обзиром на (28) -- §~4 -- показује, же $C^{(i)}(x)$ јест делимы чрез $\mideal/\gideal_{i-1}$; тут под $C^{(i)}(x)$ разумеје се којаколи не идентично изчезајуча $\rho$-рядна детерминанта из $\Mmod^*_{i-1}$. Понеже $C^{(i)}$ за $i>1$ јест нулового степена в $x_1$, модул $\mathfrak A_1$, одповедајучи идеалу $\aideal=\aideal_1$, содержи полиномы $C^{(i)},x_1C^{(i)},\ldots,x_1^{k_1-1}C^{(i)}$, и зато $(C^{(i)})^{k_1}$ јест делимы чрез $\aideal_2$; подобно $(C^{(i)})^{k_1k_2}$ јест делимы чрез $\aideal_3$, и наконец $(C^{(i)})^{k_1k_2\cdots k_{i-1}}$ чрез $\aideal_i$. Тако $\aideal_1,\ldots,\aideal_i$ сут различно од нулового идеала, што важи такоже за $i=1$. Нехай ныне $D^{(i)}(x)$ буде највечши обчи делитель всих полиномов из $\aideal_i$, значи толико тогда од степена $p_i>0$, когда $\aideal_{i+1}$ ставаје једнак нуловому идеалу. По дефиницији $D^{(i)}(x)$ и с обзиром на делительност $\aideal_i$ чрез $\aideal$ добивамо:
\[
 d^{(i+1)}D^{(i)}(x)\equiv0(\mideal/\gideal_{i-1});\qquad d^{(i+1)}\ne0;
\]
зато $d^{(i+1)}D^{(i)}$ ставаје полиномом из $\mideal/\gideal_{i-1}$, свободным од $x_1\ldots x_{i-1}$. Но $E^{(i)}(x)$ был дефинованы како највышши элементарны делитель $\Mmod_{i-1}$ односно $\Mmod^*_{i-1}$ (Теорем II и §~4) како највечши обчи делитель -- в полиномном смыслу -- всих полиномов из $\mideal/\gideal_{i-1}$, свободных од $x_1\ldots x_{i-1}$. Зато $d^{(i+1)}D^{(i)}$ и, ради регуларности $E^{(i)}(x)$ взходно к $x_i$, такоже $D^{(i)}(x)$ јест делимы чрез $E^{(i)}(x)$.

То само разсудјенье можно провести, ако трансформација (12) была од почетка сложена с (35), т. ј. ако неопредельене $u_{\mu\nu}$ былы заменьене чрез $w_{\mu\nu}$; то дава делительност $H^{(i)}(z,x)$ чрез одповедајучи $\bar E^{(i)}(z,x)$. Понеже, подобно како $D^{(i)}(x)$ и $H^{(i)}(z,x)$ једно из другого наставају специализацијеју, такоже $E^{(i)}(x)$ и $\bar E^{(i)}(z,x)$, и исто $R^{(i)}(x)$ и $\bar R^{(i)}(z,x)$, меджусобно наставају специализацијеју, морају и числа многократности при разложеньу на линеарне факторы меджусобно совпадати. Ако се јеште взаме в обзир, же некоја потяга $E^{(i)}(x)$ јест делима чрез $R^{(i)}(x)$, добива се в суме:

\medskip
\noindent\textbf{Теорем XIII.} {\itshape Ако $R^{(i)}(x)$ разложи се в алгебраичном разширительнем польу од $P(u,x_{i+1},\ldots,x_n)$ на линеарне факторы $x_i-\bar x_i$, тогда $\bar x_i$ можно на једен и само једен способ дополнити до приналежного система величин $\bar x_1\ldots\bar x_i,x_{i+1}\ldots x_n$, кторы јест нула $\mideal/\gideal_{i-1}$ и следовательно нула $\mideal$. Конечно многе нулы $\mideal$, кторе тако наставају из линеарных факторов $R^{(i)}$ -- конечно многе по присојединьеньу $x_{i+1}\ldots x_n$ -- собирајут се чрез експлицитно разложенье:}
\begin{equation}
 \bar R^{(i)}(z,x)=\prod\{z_i-(\bar x_i+v_1\bar x_1+\cdots+v_{i-1}\bar x_{i-1})\}^{\alpha_\lambda}.
\tag{36}
\end{equation}
{\itshape То разложенье, ради регуларности $\bar R^{(i)}(z,x)$ взходно к $z_i$, оставаје сохраньено, ако $x_{i+1}\ldots x_n$, присојединьене к $P(u)$, заменити којимиколи специалными системами величин $\bar x_{i+1}\ldots\bar x_n$ из $P(u)$ или из алгебраично затворјеного польа, выведеного из ньега.}

Тым самым проведено јест такоже раздельенье нул $\mideal$ по појединым факторам резултантној формы; число неопредельених $x$, присојединьених к $P(u)$, называ се такоже димензија одповедајучего алгебраичного образованьа.

Ако в разложеньу $\bar R^{(i)}(z,x)$ или в одповедајучем разложеньу $R^{(i)}(x)$ собрати факторы конјуговане взходно к $P(u,x_{i+1}\ldots x_n)$, кторе при обчем $P$ могут быти целком или честично идентичне, добива се разложенье $R^{(i)}(x)$ на потягы полиномов неразложимых взходно к $P(u,x_{i+1}\ldots x_n)$:
\[
 R^{(i)}(x)=S^{(i)}_1(x)^{\beta_1}\cdots S^{(i)}_\nu(x)^{\beta_\nu}.
\]
Тому разложеньу по Теорему X одповедаје представјенье $\gideal_i=[\jideal_1\ldots\jideal_\nu]$ тако, же $S^{(i)}_\mu(x)^{\beta_\mu}$ јест еднако норме $\gideal_{i-1}$ взходно к идеалу $\jideal_\mu$, односно норме модула $\Gmod_{i-1}$ взходно к модулу, одповедајучему $\jideal_\mu$. Тым самым објасньено јест такоже значенье експонентов $\beta_\mu$; особливо, ако $\gamma_\mu$ означује степень $S^{(i)}_\mu$, тогда многократност $\beta_\mu\gamma_\mu$ јест једнака числу остатковых класов $\gideal_{i-1}$ модуло $\jideal_\mu$, линеарно независных взходно к $P(u,x_{i+1}\ldots x_n)$.

Кратко разсмотримо јеште специалны случай, где $P$ јест карактеристики нул. Ако тут специализовати $u_{\mu\nu}$ до величин $\bar u_{\mu\nu}$ из $P$ тако, же $n$ детерминантов $C^{(i)}(x)$ $(i=1,2,\ldots,n)$ всагда оставајут регуларне взходно к $x_i$,\footnote{\mbox{Hentzelt} од самого почетка дела ту специализацију за $u_{\mu\nu}$ и зато при доказу разложимости $H^{(i)}(z,x)$ односно $\bar R^{(i)}(z,x)$ на линеарне факторы в $z$ и $v$ муси привлачити много сложенејше разсудјеньа. Експоненты, кторе при том наставају, не мусет необходно совпадати с вышеј даными. (Е. Н.)} тогда регуларност $R^{(i)}$ такоже сохраньа се. При тој специализацији $R^{(i)}$ и такисто $\bar R^{(i)}(z,x)$ ставајут обчи делитель всих $\rho$-рядных детерминантов из $\Mmod^*_{i-1}$, але не необходно највечши обчи делитель. Но специализацију можно всагда избрати тако, же и то последнье свойство оставаје сохраньено. Бо $R^{(i)}(z,x)$ можно -- како показује ексистенција базиса модула -- такоже карактеризовати како највечши обчи делитель конечно многих $\rho$-рядных детерминантов из $\Mmod^*_{i-1}$. Разложенье тако специализованого $\bar R^{(i)}(z,x)$ тогда добива се просто заменоју $u_{\mu\nu}$ чрез $\bar u_{\mu\nu}$ в (36), тако же цела теория елиминације оставаје сохраньена.

\begin{center}
(Доставльено 17.3.1922.)
\end{center}

\end{document}

